% ============================================================
% پایه‌های الگوریتم‌های کوانتومی و شبیه‌سازی
% سند مطالعاتی زنده: قطب‌نما، نه کتاب درسی
% ============================================================

\chapter*{این سند چگونه خوانده شود}
\addcontentsline{toc}{chapter}{این سند چگونه خوانده شود}

مخاطب این متن خودِ پژوهشگر است. هدف آن نگه‌داشتن شمای کلی است تا در پیچ‌وخم مدار و فرمول گم نشوید. از دیگر سو برای روش ها و رهیافت های اصلی توضیحات و توصیفات مفصل اورده خواهد شد تا بخش مربوطه مرجعی باشد برای فهم ریاضیات و فیزیک مسئله. هر فصل کوتاه است، شهودی است، و همان فایل بعداً گسترش می‌یابد.

سه لایه را از هم جدا نگه دارید؛ درهم‌ریختگی معمولاً از قاطی کردن این‌هاست:
\begin{enumerate}
	\item \textbf{هدف / مسئله} --- چه چیزی را می‌خواهیم بدانیم؟
	\item \textbf{رهیافت محاسباتی} --- با چه مدل محاسباتی به آن می‌رسیم؟
	\item \textbf{ابزار پیاده‌سازی} --- کدام زیرروال آن رهیافت را محقق می‌کند؟
\end{enumerate}
نقشهٔ زنده در فصل~\ref{chap:map} است. اگر وسط کار پژوهشی سردرگم شدید، به همان فصل و به قطب‌نمای فصل~\ref{chap:compass} برگردید.

منبع واحد واژگان جدول~\ref{tab:glossary} در فصل~\ref{chap:glossary} است. در متن روان، فارسی مصوب می‌آید و شکل انگلیسی در نخستین کاربرد هر فصل داخل پرانتز است. سرواژه‌ها و نام الگوریتم‌ها (\lr{Trotter}، \lr{LCU}، \lr{QSP}، \lr{QPE}، \lr{VQE}) انگلیسی می‌مانند.

نسخهٔ اول عمداً اثبات پیچیدگی، جزئیات مدار، و کد شبیه‌ساز ندارد؛ این‌ها در «جای توسعه» هر فصل فهرست شده‌اند.

% ============================================================
\part{نقشهٔ مفهومی}
% ============================================================

\chapter{نقشهٔ زنده}
\label{chap:map}

\section{جایگاه در نقشه}

این فصل خودِ نقشه است. بقیهٔ فصل‌ها به اینجا ارجاع می‌دهند.

\section{ایدهٔ شهودی}

پژوهش شبیه‌سازی کوانتومی\footnote{\lr{quantum simulation}.} یک خط مستقیم از «فیزیک» به «مدار» نیست. نخست یک مسئلهٔ فیزیکی به مدل ریاضی و سپس به هامیلتونی $H$ تبدیل می‌شود. بعد باید دو چیز را جدا کرد:
\begin{itemize}
	\item \textbf{چه پرسشی از $H$ دارید؟} تحول زمانی، یا طیف / حالت پایه؟
	\item \textbf{با کدام رهیافت به آن می‌رسید؟} پیاده‌سازی $U(t)$ یا $f(H)$، مسیر آدیاباتیک، مهندسی آنالوگ، یا حلقهٔ وردشی.
\end{itemize}
ابزارهایی مثل \lr{Trotter}، \lr{LCU} و \lr{QSP} زیرِ رهیافت‌اند، نه هم‌سطح آن. برآورد فاز\footnote{\lr{phase estimation} (\lr{QPE}).} روی مرز دینامیک و طیف می‌نشیند، چون به تحول کنترل‌شده نیاز دارد.

اگر وسط جزئیات گم شدید، به این فصل برگردید و بپرسید: سؤال من کدام است؟ رهیافت کدام است؟ ابزار کدام است؟

\section{حداقل فرمالیسم}

جریان کاری\footnote{\lr{workflow}.} کلی در شکل~\ref{fig:research-map} آمده است.

\begin{figure}[htbp]
	\centering
	\begin{latin}
		\setstretch{1}
		\resizebox{\ifdim\width>\textwidth\textwidth\else\width\fi}{!}{%
			% نقشهٔ پژوهشی: لایه‌های مسئله / پرسش / رهیافت / ابزار / اجرا
% اتصالات مطابق فصل نقشه: QPE از تحول تغذیه می‌شود و پرسش طیفی را جواب می‌دهد؛
% آنالوگ از ابزار گیتی عبور نمی‌کند؛ آدیاباتیک رهیافت است نه ورودی QPE.
% سبک مشترک نمودارهای سند پایه‌ها (فقط داخل محیط latin فراخوانی شود)
\tikzset{
	qafbox/.style={
		draw=black!55,
		rounded corners=2pt,
		align=center,
		inner sep=3pt,
		thick,
		fill=white,
		font=\footnotesize\linespread{1}\selectfont,
		minimum height=0.95cm
	},
	qafwide/.style={
		qafbox,
		text width=2.55cm
	},
	qafnarr/.style={
		qafbox,
		text width=2.05cm
	},
	qaftitle/.style={
		font=\small\bfseries\linespread{1}\selectfont,
		align=center
	},
	qafarr/.style={
		-{Latex[length=2.2mm]},
		thick,
		draw=black!65
	},
	qafpanel/.style={
		draw,
		thick,
		rounded corners=4pt,
		inner xsep=8pt,
		inner ysep=8pt
	}
}

\begin{tikzpicture}[
	font=\scriptsize\linespread{1}\selectfont,
	x=1cm,
	y=1cm
	]
	\tikzset{
		lab/.style={
			font=\tiny\bfseries\linespread{1}\selectfont,
			anchor=east,
			inner sep=1pt,
			text=black!55,
			align=right
		},
		note/.style={
			font=\tiny\linespread{1}\selectfont,
			inner sep=1pt,
			text=black!60
		}
	}

	% ----- لایه مسئله -----
	\node[qafwide, fill=blue!12] (phys) at (0,0) {Physical problem};
	\node[qafwide, fill=blue!10] (model) at (0,-1.25) {Mathematical model};
	\node[qafwide, fill=blue!14] (H) at (0,-2.5) {Hamiltonian $H$};

	% ----- لایه پرسش -----
	\node[qafnarr, fill=green!14] (dyn) at (-3.05,-4.05) {Dynamic\\question};
	\node[qafnarr, fill=orange!14] (stat) at (3.05,-4.05) {Static / spectral\\question};

	% ----- لایه رهیافت -----
	\node[qafnarr, fill=cyan!14] (ana) at (-5.15,-5.7) {Analog\\engineering};
	\node[qafnarr, fill=green!10] (evo) at (-2.15,-5.7) {Time evolution\\$f(H)$};
	\node[qafnarr, fill=orange!10] (adi) at (2.15,-5.7) {Adiabatic path\\$H(s)$};
	\node[qafnarr, fill=yellow!18] (var) at (5.15,-5.7) {Variational\\hybrid};

	% ----- لایه ابزار -----
	\node[qafbox, text width=1.45cm, fill=green!6] (trot) at (-3.35,-7.35) {Trotter};
	\node[qafbox, text width=1.45cm, fill=green!6] (lcu) at (-2.15,-7.35) {LCU};
	\node[qafbox, text width=1.45cm, fill=green!6] (qsp) at (-0.95,-7.35) {QSP / QSVT};
	\node[qafbox, text width=1.7cm, fill=orange!12] (qpe) at (1.15,-7.35) {QPE};
	\node[qafbox, text width=1.55cm, fill=yellow!12] (vqe) at (5.15,-7.35) {VQE};

	% ----- لایه اجرا و خروجی -----
	\node[qafwide, fill=gray!14] (exec) at (0,-9.15) {Execution:\\circuit or analog schedule};
	\node[qafwide, fill=gray!8] (meas) at (0,-10.4) {Measurement / observable};
	\node[qafwide, fill=gray!8] (val) at (0,-11.55) {Classical validation};
	\node[qafwide, fill=gray!8] (res) at (0,-12.7) {Resource estimation};

	% برچسب لایه
	\node[lab] at (-6.55,-4.05) {QUESTION};
	\node[lab] at (-6.55,-5.7) {APPROACH};
	\node[lab] at (-6.55,-7.35) {TOOL};
	\node[lab] at (-6.55,-9.15) {RUN};

	% پس‌زمینه دو ستون پرسش
	\begin{scope}[on background layer]
		\node[qafpanel, fill=green!5, draw=green!40!black, fit=(dyn)(evo)(trot)(lcu)(qsp),
			inner xsep=5pt, inner ysep=6pt] {};
		\node[qafpanel, fill=orange!5, draw=orange!45!black, fit=(stat)(adi)(var)(vqe),
			inner xsep=5pt, inner ysep=6pt] {};
	\end{scope}

	% ----- اتصالات مسئله -----
	\draw[qafarr] (phys) -- (model);
	\draw[qafarr] (model) -- (H);
	\coordinate (hsplit) at (0,-3.15);
	\draw[qafarr] (H.south) -- (hsplit) -| (dyn.north);
	\draw[qafarr] (hsplit) -| (stat.north);

	% پرسش → رهیافت
	\draw[qafarr] (dyn) -- (evo);
	\draw[qafarr] (dyn.west) -- ++(-0.35,0) |- (ana.north);
	\draw[qafarr] (stat) -- (adi);
	\draw[qafarr] (stat.east) -- ++(0.35,0) |- (var.north);
	% پرسش طیفی می‌تواند از رهیافت تحول هم استفاده کند (QPE یا $f(H)$)
	\draw[qafarr, dashed] (stat.west) -- (evo.east)
		node[note, midway, above] {via $U(t)$};

	% رهیافت تحول → سه ابزار دیجیتال
	\draw[qafarr] (evo.south) -- ++(0,-0.18) -| (trot.north);
	\draw[qafarr] (evo) -- (lcu);
	\draw[qafarr] (evo.south) -- ++(0,-0.18) -| (qsp.north);

	% وردشی → VQE
	\draw[qafarr] (var) -- (vqe);

	% QPE از هر پیاده‌سازی U تغذیه می‌شود، نه فقط از یک ابزار
	\node[draw=green!45!black, densely dashed, rounded corners=2pt, inner sep=3pt,
		fit=(trot)(lcu)(qsp)] (utools) {};
	\draw[qafarr] (utools.east) -- (qpe.west)
		node[note, above, pos=0.55] {controlled $U$};

	% شینهٔ اجرا: همهٔ مسیرها جداگانه به اجرا می‌رسند؛ آنالوگ از گیت‌ها عبور نمی‌کند
	\coordinate (busL) at (-5.15,-8.35);
	\coordinate (busR) at (5.15,-8.35);
	\draw[thick, draw=black!55] (busL) -- (busR);
	\foreach \n in {ana,trot,lcu,qsp,qpe,adi,vqe} {
		\draw[qafarr] (\n.south) -- (\n.south |- busL);
	}
	\draw[qafarr] (0,-8.35) -- (exec.north);

	\draw[qafarr] (exec) -- (meas);
	\draw[qafarr] (meas) -- (val);
	\draw[qafarr] (val) -- (res);
\end{tikzpicture}
%
		}
	\end{latin}
	\caption{نقشهٔ پژوهشی در چهار لایهٔ پس از $H$: پرسش، رهیافت، ابزار، اجرا. کادر خط‌چین دور \lr{Trotter}/\lr{LCU}/\lr{QSP} یعنی هر سه پیاده‌سازی $U$ هستند و \lr{QPE} از همان گروه تغذیه می‌شود. پیکان خط‌چین از پرسش طیفی به \lr{QPE} نشان می‌دهد این ابزار پرسش طیفی را جواب می‌دهد ولی از لایهٔ رهیافت تحول عبور می‌کند. مسیر آدیاباتیک ابزار جداگانه ندارد و آنالوگ بدون گیت دیجیتال به اجرا می‌رسد.}
	\label{fig:research-map}
\end{figure}

اتصالات شکل~\ref{fig:research-map} باید با همین تفکیک خوانده شوند:
\begin{itemize}
	\item \lr{Trotter} / \lr{LCU} / \lr{QSP} ابزارِ رهیافت تحول / $f(H)$ هستند؛ $U(t)=e^{-iHt}$ فقط یک کاربرد است. هر سه، جدا از هم، به اجرا می‌رسند.
	\item \lr{QPE} ابزار طیفی است اما ورودی‌اش تحول کنترل‌شده است؛ بنابراین از جعبهٔ ابزارهای $U$ تغذیه می‌شود، نه از مسیر آدیاباتیک و نه فقط از \lr{LCU}.
	\item مسیر آدیاباتیک $H(s)$ یک \emph{رهیافت} موازی با تحول و وردشی است، نه زیرروال \lr{Trotter} و نه پیش‌نیاز \lr{QPE}.
	\item \lr{VQE} زیرِ رهیافت وردشی / ترکیبی است، نه تجزیهٔ $e^{-iHt}$.
	\item مهندسی آنالوگ از پرسش دینامیکی می‌آید و بدون عبور از گیت‌های دیجیتال به اجرا می‌رسد.
\end{itemize}


\begin{table}[htbp]
	\centering
	\footnotesize
	\setstretch{1.15}
	\setlength{\tabcolsep}{4pt}
	\caption{از پرسش علمی به رهیافت و ابزار نوعی.}
	\label{tab:question-approach-tool}
	\begin{tabularx}{\textwidth}{@{}
		>{\raggedright\arraybackslash}p{4.8cm}
		>{\raggedright\arraybackslash}p{2.5cm}
		>{\raggedright\arraybackslash}p{3.9cm}
		>{\raggedright\arraybackslash}p{3.5cm}
%		>{\raggedright\arraybackslash}p{1.1cm}
		@{}}
		\toprule
		\rowcolor{gray!14}
		\textbf{اگر می‌خواهید} & \textbf{پرسش} & \textbf{رهیافت} & \textbf{ابزار نوعی} 
%		& \textbf{فصل} 
		\\
		\midrule
		$\lvert\psi(t)\rangle$ یا مشاهده‌پذیر وابسته به \lr{t}
		& دینامیک
		& تحول زمانی / $f(H)$
		& \lr{Trotter}، \lr{LCU}، \lr{QSP}
%		& \ref{chap:dyn-spec}، \ref{chap:digital-tools}
		\\
		\rowcolor{gray!8}
		انرژی پایه یا گاف
		& ایستا / طیفی
		& وردشی، آدیاباتیک، یا \lr{QPE}
		& \lr{VQE}، مسیر $H(s)$، \lr{QPE}
%		& \ref{chap:dyn-spec}، \ref{chap:evo-adi}، \ref{chap:spectral-tools}
		\\
		ویژه مقدارهای $H$ با دقت بالا
		& طیفی
		& تحول کنترل‌شده $+$ برآورد فاز
		& شبیه‌سازی $U$ سپس \lr{QPE}
%		& \ref{chap:spectral-tools}
		\\
		\rowcolor{gray!8}
		تابع پاسخ / \lr{Green}
		& پل دینامیک--طیف
		& $f(H)$ یا تحول
		& \lr{QSP}، \lr{QPE}
%		& \ref{chap:dyn-spec}
		\\
		مهندسی مستقیم $H$ روی سخت‌افزار
		& دینامیک یا ایستا
		& آنالوگ
		& نگاشت $H_{\mathrm{target}}$
%		& \ref{chap:hardware}
		\\
		\rowcolor{gray!8}
		ماندن در حالت پایه لحظه‌ای
		& ایستا
		& آدیاباتیک
		& زمان‌بندی $H(s)$
%		& \ref{chap:evo-adi}
		\\
		\bottomrule
	\end{tabularx}
\end{table}

\section{ارتباط}

قرارداد واژگان در فصل~\ref{chap:glossary} است. سه لایهٔ مفهومی در فصل~\ref{chap:three-layers}، جریان کاری با مثال در فصل~\ref{chap:pipeline}، و درخت تصمیم در فصل~\ref{chap:compass} آمده است.

\section{چه وقت استفاده می‌شود}

هر بار که یک مقاله، یک روش، یا یک ایدهٔ پژوهشی جدید دیدید و ندانستید زیر کدام شاخه است.

\begin{remark}[جای توسعه]
افزودن سطر برای مسائل گرمایی، سامانهٔ کوانتومی باز، و ترابرد؛ تفکیک \lr{digital adiabatic} از آدیاباتیک آنالوگ در خودِ نمودار.
\end{remark}

% ============================================================
\part{پایه‌های محاسباتی}
% ============================================================

\chapter{الگوریتم چیست؟}
\label{chap:algorithm}

\section{جایگاه در نقشه}

لایهٔ صفر: قبل از کوانتوم. بدون این تعریف، «الگوریتم کوانتومی» با سخت‌افزار یا با یک گیت خاص قاطی می‌شود. در نقشهٔ فصل~\ref{chap:map} هنوز به هامیلتونی نرسیده‌ایم؛ اینجا فقط می‌گوییم \emph{دستور حل مسئله} چیست.

\section{ایدهٔ شهودی}

الگوریتم یک دستورالعمل محاسباتی است: ورودی مشخص می‌گیرد، با تعداد متناهی گامِ بدون ابهام پیش می‌رود، و خروجی می‌دهد. رایانه، آزمایشگاه، یا کاغذ و قلم مجری‌اند، نه خودِ الگوریتم.

سه پرسش همیشه همراه الگوریتم‌اند: \textbf{درستی} (آیا خروجی همان چیزی است که مسئله می‌خواهد؟)، \textbf{پایان‌پذیری} (آیا حتماً تمام می‌شود؟)، و \textbf{هزینه} (با اندازهٔ ورودی، زمان و حافظه چطور رشد می‌کنند؟).

شهود مفید: دستور پخت غذا الگوریتم نیست اگر «کمی نمک» مبهم بماند. الگوریتم وقتی الگوریتم است که اجراکننده‌های مختلف، با همان ورودی، به همان نتیجه برسند (تا حد خطای اعلام‌شده).

\section{حداقل فرمالیسم}

یک الگوریتم $A$ را می‌توان نگاشتی از نمونهٔ مسئله $x$ به خروجی $A(x)$ دانست، به‌همراه مدلی از هزینه (تعداد گام، عمق مدار، تعداد پرسش از یک اوراکل). در این سند دو تمایز کافی است:
\begin{itemize}
	\item \textbf{مسئله} --- چه چیزی پرسیده می‌شود؟ (مثلاً انرژی پایهٔ یک مدل)
	\item \textbf{الگوریتم} --- با چه روالی به پاسخ می‌رسیم؟
\end{itemize}
پیچیدگی را اینجا وارد جزئیات نمی‌کنیم. «کوانتومی بودن» مدل اجرا را عوض می‌کند، نه تعریف الگوریتم را.

\section{ارتباط}

خانواده‌های کلاسیک مرتبط در فصل~\ref{chap:classical}، مدل محاسباتی کوانتومی در فصل~\ref{chap:qmodel}، و تمایز الگوریتم در برابر شبیه‌سازی در برابر ابزار در فصل~\ref{chap:three-layers} آمده است.

\section{چه وقت استفاده می‌شود}

وقتی وسوسه می‌شوید \lr{Trotter} یا \lr{VQE} را «خودِ مسئله» بنامید. آن‌ها الگوریتم یا زیرروال‌اند؛ مسئله چیز دیگری است (مثلاً تحول یک مدل هابارد).

\begin{remark}[جای توسعه]
مدل‌های هزینه (مدار، پرسمان، فضا)؛ تمایز الگوریتم دقیق، تقریبی، تصادفی و اکتشافی؛ اشاره به Church--Turing و نسخهٔ کوانتومی آن.
\end{remark}

\chapter{خانواده‌های کلاسیکِ مرتبط}
\label{chap:classical}

\section{جایگاه در نقشه}

هنوز در لایهٔ کلاسیک هستیم. هدف، کاتالوگ الگوریتم‌های عمومی نیست؛ فقط خانواده‌هایی که بعداً با شبیه‌سازی یا الگوریتم کوانتومی تلاقی مفهومی دارند.

\section{ایدهٔ شهودی}

هر بار که یک الگوریتم کوانتومی می‌بینید، معمولاً جانشین یا رقیب یک روال کلاسیک است. اگر آن رقیب را نشناسید، نمی‌فهمید کوانتوم دقیقاً چه چیزی را عوض کرده: مدل هزینه را، نمایش حالت را، یا خودِ سؤال را. پنج خانواده برای این سند کافی‌اند.

\paragraph{جستجو و بهینه‌سازی.}
جستجوی ناساخت‌یافته و بهینه‌سازی ترکیبیاتی. در کوانتوم: \lr{Grover}، \lr{QAOA}، و بازپخت کوانتومی.\footnote{\lr{quantum annealing}.}

\paragraph{تبدیل فوریه.}
از حوزهٔ زمان/مکان به بسامد. نسخهٔ کوانتومی آن \lr{QFT} است و ستون برآورد فاز.

\paragraph{دستگاه خطی و زیرفضای \lr{Krylov}.}
حل $Ax=b$ و استخراج اطلاعات طیفی با تکرارهای ماتریسی. در کوانتوم: \lr{HHL} و به‌طور کلی‌تر توابعی از عملگر از راه بلوک‌کدگذاری.\footnote{\lr{block-encoding}.} شبیه‌سازی هامیلتونی هم از همین جنس است: اعمال تابعی از $H$.

\paragraph{نمونه‌برداری و مونت‌کارلو.}
برآورد یک مقدار با نمونه‌های تصادفی. نسخهٔ کوانتومی اغلب برآورد دامنه است. در فیزیک بس‌ذره‌ای،\footnote{\lr{many-body}.} مونت‌کارلو کوانتومی خودش یک رقیب کلاسیک جدی است.

\paragraph{شبیه‌سازی کلاسیک سامانه‌های کوانتومی.}
این مهم‌ترین رقیب مستقیم پژوهش شماست: قطری‌سازی دقیق و شبیه‌سازی بردار حالت\footnote{\lr{state vector simulation}.} (دقیق، نمایی در تعداد ذرات)؛ مونت‌کارلو کوانتومی (گاهی مقیاس بهتر، با محدودیت علامت)؛ و شبکهٔ تانسوری\footnote{\lr{tensor network}.} و حالت حاصل‌ضرب ماتریسی\footnote{\lr{matrix product state}.} از جمله \lr{DMRG}، که برای درهم‌تنیدگی محدود قوی‌اند~\cite{Verstraete2008,Cirac2021}. الگوریتم کوانتومی وقتی معنا دارد که یکی از این مسیرها از نظر مقیاس، دقت، یا نوع سؤال ناکافی شود.

\section{حداقل فرمالیسم}

هزینهٔ کلاسیک شبیه‌سازی یک سامانهٔ کوانتومی عمومی معمولاً با بُعد فضای هیلبرت --- یعنی به‌طور نمایی با تعداد ذرات --- رشد می‌کند. استثناها همان روش‌هایی‌اند که ساختار (محلی بودن، نبود علامت، درهم‌تنیدگی کم) را استثمار می‌کنند.

\section{ارتباط}

تعریف الگوریتم در فصل~\ref{chap:algorithm}، مدل کوانتومی در فصل~\ref{chap:qmodel}، مدل‌های فیزیکی در فصل~\ref{chap:models}، و راستی‌آزمایی با حد کلاسیک در فصل~\ref{chap:closing} آمده است.

\section{چه وقت استفاده می‌شود}

قبل از انتخاب یک الگوریتم کوانتومی بپرسید: همین سؤال را کلاسیک تا چه اندازه و با چه روشی می‌توانم جواب بدهم؟

\begin{remark}[جای توسعه]
جدول مقیاس‌بندی کیفی میان قطری‌سازی، \lr{MPS}، \lr{QMC} و مدار کوانتومی؛ قانون مقیاس‌بندی هزینه در هر خانواده.
\end{remark}

\chapter{مدل محاسباتی کوانتومی}
\label{chap:qmodel}

\section{جایگاه در نقشه}

این فصل می‌گوید رایانهٔ کوانتومی و الگوریتم کوانتومی \emph{کجا} روی نقشه می‌نشینند: یک مدل اجرای تازه برای همان مفهوم الگوریتم، نه صرفاً رایانه‌ای سریع‌تر.

\section{ایدهٔ شهودی}

رایانهٔ کلاسیک بیت را در $0$ یا $1$ نگه می‌دارد. رایانهٔ کوانتومی با حالت‌هایی در فضای هیلبرت کار می‌کند. سه منبع محاسبات را از کلاسیک جدا می‌کنند: برهم‌نهی،\footnote{\lr{superposition}.} تداخل،\footnote{\lr{interference}.} و درهم‌تنیدگی.\footnote{\lr{entanglement}.}

الگوریتم کوانتومی روالی است که از این منابع برای یک مسئله با ورودی کلاسیک یا کوانتومی استفاده می‌کند و در پایان با اندازه‌گیری به خروجی کلاسیک (یا به یک حالت کوانتومی مفید) می‌رسد~\cite{NielsenChuang2010}. شبیه‌سازی کوانتومی یک کلاس از مسئله‌ها روی همین مدل است، نه تعریف خودِ مدل~\cite{Feynman1982,Lloyd1996,Georgescu2014}.

سخت‌افزار سه شکل اصلی دارد که بعداً تکرار می‌شوند:
\begin{enumerate}
	\item \textbf{مدار گیتی (دیجیتال):} دنبالهٔ گیت‌ها از یک مجموعه گیت جهانی؛\footnote{\lr{universal gate set}.}
	\item \textbf{آنالوگ:} مهندسی خودِ هامیلتونی روی دستگاه؛
	\item \textbf{آدیاباتیک:} محاسبه با مسیر کند $H(s)$.
\end{enumerate}
این سه، سه رهیافت اجرا هستند. یک هدف پژوهشی (مثلاً انرژی پایهٔ هابارد) ممکن است با هر سه دنبال شود.

\section{حداقل فرمالیسم}

مدل مدار: حالت $\lvert\psi\rangle$ روی $n$ کیوبیت، تحول یکانی $U$ ساخته‌شده از گیت‌های محلی، سپس اندازه‌گیری در یک پایه. خروجی الگوریتم نوعاً یک برآورد از $\langle\psi\rvert O\lvert\psi\rangle$ است، نه کل تابع موج مگر در مسائل کوچک.

تمایز کوتاه: گیت\footnote{\lr{gate}.} دستور موضعی است؛ مدار الگوریتم در مدل دیجیتال است؛ زمان‌بندی هامیلتونی الگوریتم در مدل آنالوگ یا آدیاباتیک است.

\section{ارتباط}

الگوریتم به‌معنای عام در فصل~\ref{chap:algorithm}، سه لایه در فصل~\ref{chap:three-layers}، دیجیتال در برابر آنالوگ در فصل~\ref{chap:hardware}، و تحول در برابر آدیاباتیک در فصل~\ref{chap:evo-adi} آمده است.

\section{چه وقت استفاده می‌شود}

وقتی می‌پرسید الگوریتم‌های کوانتومی چه هستند، و شبیه‌سازی و رایانهٔ کوانتومی کجای این بحث‌اند. پاسخ کوتاه: رایانه مدل اجراست؛ الگوریتم روال است؛ شبیه‌سازی یک خانواده از مسائل روی همان مدل است.

\begin{remark}[جای توسعه]
مدل پرسمان و \lr{BQP} در حد یک پاراگراف؛ نویز و ناهمدوسی به‌عنوان فاصلهٔ مدل آرمانی از دستگاه؛ \lr{early-FTQC} در برابر \lr{NISQ}.
\end{remark}

\chapter{سه لایه: الگوریتم، شبیه‌سازی، ابزار}
\label{chap:three-layers}

\section{جایگاه در نقشه}

این فصل ستون مفهومی کل سند است. درهم‌ریختگی پژوهشی معمولاً از قاطی کردن این سه لایه می‌آید. جدول عملی «اگر می‌خواهید\ldots» در فصل~\ref{chap:map} است.

\section{ایدهٔ شهودی}

سه مفهوم نزدیک‌اند ولی یکی نیستند. شکل~\ref{fig:three-layers} ترتیب منطقی‌شان را نشان می‌دهد.

\begin{figure}[htbp]
	\centering
	\begin{latin}
		\setstretch{1}
		% سبک مشترک نمودارهای سند پایه‌ها (فقط داخل محیط latin فراخوانی شود)
\tikzset{
	qafbox/.style={
		draw=black!55,
		rounded corners=2pt,
		align=center,
		inner sep=3pt,
		thick,
		fill=white,
		font=\footnotesize\linespread{1}\selectfont,
		minimum height=0.95cm
	},
	qafwide/.style={
		qafbox,
		text width=2.55cm
	},
	qafnarr/.style={
		qafbox,
		text width=2.05cm
	},
	qaftitle/.style={
		font=\small\bfseries\linespread{1}\selectfont,
		align=center
	},
	qafarr/.style={
		-{Latex[length=2.2mm]},
		thick,
		draw=black!65
	},
	qafpanel/.style={
		draw,
		thick,
		rounded corners=4pt,
		inner xsep=8pt,
		inner ysep=8pt
	}
}

\begin{tikzpicture}[
	font=\footnotesize\linespread{1}\selectfont,
	node distance=0.55cm
	]
	\node[qafbox, text width=7.6cm, fill=blue!10] (l2) {%
		\textbf{Layer 2 --- Quantum simulation}\\[2pt]
		Which physical model is being represented?};
	\node[qafbox, text width=7.6cm, fill=green!10, below=of l2] (l1) {%
		\textbf{Layer 1 --- Quantum algorithm}\\[2pt]
		Which procedure solves that question?};
	\node[qafbox, text width=7.6cm, fill=orange!10, below=of l1] (l3) {%
		\textbf{Layer 3 --- Implementation tools}\\[2pt]
		Trotter, LCU, QSP, QPE, VQE, adiabatic path};
	\draw[qafarr] (l2) -- (l1);
	\draw[qafarr] (l1) -- (l3);
\end{tikzpicture}

	\end{latin}
	\caption{سه لایهٔ نزدیک ولی جدا: شبیه‌سازی می‌گوید چه مدلی مطالعه می‌شود؛ الگوریتم می‌گوید با چه روالی؛ ابزار می‌گوید آن روال چطور ساخته می‌شود.}
	\label{fig:three-layers}
\end{figure}

\paragraph{لایهٔ ۱ --- الگوریتم کوانتومی.}
دستورالعمل محاسباتی برای حل یک مسئله. مثال: با این مدار، انرژی پایه را تا خطای $\varepsilon$ برآورد کن.

\paragraph{لایهٔ ۲ --- شبیه‌سازی کوانتومی.}
استفاده از یک سامانهٔ کوانتومی برای بازنمایی و مطالعهٔ یک سامانه یا مدل کوانتومی دیگر. اینجا هدف پژوهشی است: مدل هابارد، یک مولکول، یک نظریهٔ پیمانه‌ای شبکه‌ای.\footnote{\lr{lattice gauge theory}.} شبیه‌سازی می‌تواند دیجیتال، آنالوگ، یا ترکیبی باشد. شبیه‌سازی یک گیت خاص نیست~\cite{CiracZoller2012,Georgescu2014}.

\paragraph{لایهٔ ۳ --- روش‌های پیاده‌سازی.}
ابزارهایی برای تحقق یک هدف محاسباتی مشخص: \lr{Trotter}، ترکیب خطی یکانی‌ها (\lr{LCU})، پردازش سیگنال کوانتومی (\lr{QSP})، برآورد فاز (\lr{QPE})، \lr{VQE}، و تحول آدیاباتیک.

جملهٔ جهت‌یاب: شبیه‌سازی می‌گوید چه مدلی را مطالعه می‌کنیم؛ الگوریتم می‌گوید با چه روالی؛ ابزار می‌گوید آن روال را چطور روی دستگاه می‌سازیم.

\section{حداقل فرمالیسم}

\begin{table}[htbp]
	\centering
	\footnotesize
	\setstretch{1.15}
	\caption{نمونهٔ جداسازی سه لایه.}
	\label{tab:three-layers}
	\begin{tabularx}{\textwidth}{@{}>{\raggedright\arraybackslash}p{4.4cm} c >{\raggedright\arraybackslash}X@{}}
		\toprule
		\rowcolor{gray!14}
		\textbf{می‌خواهید} & \textbf{لایه} & \textbf{نمونه} \\
		\midrule
		مسئله را نام ببرید
		& ۲
		& دینامیک هابارد پس از تغییر ناگهانی هامیلتونی (\lr{quench})
		\\
		\rowcolor{gray!8}
		روال حل را بنویسید
		& ۱
		& آماده‌سازی $\lvert\psi_0\rangle$، اعمال $U(t)$، اندازه‌گیری $O$
		\\
		$U(t)$ را بسازید
		& ۳
		& فرمول حاصل‌ضربی یا \lr{QSP}
		\\
		\rowcolor{gray!8}
		ویژه مقدار را بخوانید
		& ۱ و ۳
		& \lr{QPE} که زیرروالش شبیه‌سازی کنترل‌شده است
		\\
		حالت پایه روی \lr{NISQ}
		& ۱ و ۳
		& \lr{VQE}: آنزات وردشی $+$ بهینه‌ساز کلاسیک
		\\
		\bottomrule
	\end{tabularx}
\end{table}

یک ابزار می‌تواند در چند الگوریتم ظاهر شود. \lr{Trotter} هم دینامیک می‌سازد، هم مسیر آدیاباتیک دیجیتال، هم ورودی \lr{QPE}.

\section{ارتباط}

مدل اجرا در فصل~\ref{chap:qmodel}، مدل‌های فیزیکی در فصل~\ref{chap:models}، رهیافت‌ها در فصل~\ref{chap:evo-adi}، و ابزارها در فصل‌های~\ref{chap:digital-tools} و~\ref{chap:spectral-tools} آمده‌اند.

\section{چه وقت استفاده می‌شود}

قبل از خواندن هر مقاله بپرسید: این کار مسئله است، الگوریتم است، یا ابزار؟ اگر هر سه را در یک جمله قاطی کردید، به این فصل برگردید.

\begin{remark}[جای توسعه]
تمایز شبیه‌سازی آنالوگ فاینمن در برابر شبیه‌سازی دیجیتال لوید؛ جایگاه نمونه‌برداری بوزونی به‌عنوان مسئله‌ای مجزا از شبیه‌سازی هامیلتونی.
\end{remark}

% ============================================================
\part{از مسئلهٔ فیزیکی تا جریان کاری}
% ============================================================

\chapter{از مسئلهٔ فیزیکی تا هامیلتونی}
\label{chap:models}

\section{جایگاه در نقشه}

بالای نمودار فصل~\ref{chap:map}: مسئلهٔ فیزیکی $\to$ مدل ریاضی $\to$ هامیلتونی $H$. قبل از انتخاب \lr{Trotter} یا \lr{VQE} باید معلوم شود $H$ چیست و جمله‌هایش چه ساختاری دارند.

\section{ایدهٔ شهودی}

فیزیک خام مستقیماً الگوریتم نمی‌شود. نخست مدلی می‌سازیم که درجات آزادی را نگه می‌دارد و بقیه را دور می‌ریزد. خروجی این مرحله تقریباً همیشه یک هامیلتونی است،
\begin{equation}
	H = \sum_j h_j,
\end{equation}
که هر $h_j$ موضعی یا کم‌وزن است. از اینجا به بعد، مسیر الگوریتمی برای مدل‌های خیلی متفاوت شبیه هم می‌شود. آنچه عوض می‌شود کدگذاری درجات آزادی روی کیوبیت (یا مد بوزونی) و تعداد و نوع جمله‌هاست. چهار خانواده برای جهت‌یابی کافی‌اند؛ هیچ‌کدام را اینجا حل نمی‌کنیم.

\paragraph{اسپین: \lr{Ising} و \lr{Heisenberg}.}
ساده‌ترین ورود. درجات آزادی اسپین روی شبکه؛ $H$ جمع جمله‌های دو جسمی مثل $Z_i Z_j$ یا $\vec{S}_i\cdot\vec{S}_j$. شهود: مغناطیس، شیشهٔ اسپینی، و بسیاری مسائل بهینه‌سازی که به \lr{Ising} نگاشته می‌شوند.

\paragraph{هابارد (فرمی و بوز).}
مدل استاندارد همبستگی قوی: ذرات روی شبکه می‌پرند ($t$) و اگر روی یک سایت باشند انرژی $U$ می‌پردازند. فرمی-هابارد الکترون‌های همبسته است؛ بوز-هابارد اتم‌های سرد در شبکهٔ نوری. پرسش‌های نوعی: عایق مات، ابرشاره، ترابرد، دینامیک پس از تغییر ناگهانی هامیلتونی.

\paragraph{ساختار الکترونی مولکولی.}
الکترون‌ها در میدان هسته‌ها. پس از انتخاب پایه و اغلب یک فضای فعال،\footnote{\lr{active space}.} $H$ به شکل فرمیونی یا پائولی درمی‌آید~\cite{McArdle2020,Bauer2020}. شیمی کوانتومی از همین‌جا به \lr{VQE} و به شبیه‌سازی تحمل‌پذیر در برابر خطا وصل می‌شود.

\paragraph{میدان روی شبکه و نظریهٔ پیمانه‌ای شبکه‌ای.}
میدان کوانتومی پیوسته را روی شبکه می‌نشانند تا درجات آزادی متناهی شود. نظریه‌های پیمانه‌ای قیدهای موضعی (قانون گاوس) دارند؛ همین قیدها کدگذاری را سخت‌تر از هابارد می‌کنند.

\section{حداقل فرمالیسم}

شکل نوعی:
\begin{align}
	H_{\mathrm{spin}}
	&= \sum_{\langle ij\rangle} J_{ij}\, \sigma_i\cdot\sigma_j + \sum_i h_i\sigma_i^z, \\
	H_{\mathrm{Hubbard}}
	&= -t\sum_{\langle ij\rangle\sigma}\bigl(c_{i\sigma}^\dagger c_{j\sigma}+\mathrm{h.c.}\bigr)
	+ U\sum_i n_{i\uparrow}n_{i\downarrow}, \\
	H_{\mathrm{mol}}
	&= \sum_{pq} h_{pq} a_p^\dagger a_q
	+ \tfrac12\sum_{pqrs} h_{pqrs} a_p^\dagger a_q^\dagger a_r a_s.
\end{align}
پس از $H$، سؤال پژوهشی به فصل~\ref{chap:dyn-spec} می‌رود نه به «نوع ماده».

\section{ارتباط}

سه لایه در فصل~\ref{chap:three-layers}، جریان کاری با مثال هابارد در فصل~\ref{chap:pipeline}، رقبا کلاسیک در فصل~\ref{chap:classical}، و کدگذاری و سخت‌افزار در فصل~\ref{chap:hardware} آمده است.

\section{چه وقت استفاده می‌شود}

وقتی مقاله یا ایدهٔ تازه‌ای دیدید، اول بپرسید مدل کدام است و $H$ چند جمله و از چه نوعی دارد؛ قبل از انتخاب الگوریتم.

\begin{remark}[جای توسعه]
نگاشت فرمیون به کیوبیت (\lr{Jordan--Wigner}، \lr{Bravyi--Kitaev})؛ مدل‌های ارتعاشی--الکترونی و مد بوزونی؛ جزئیات قید پیمانه‌ای.
\end{remark}

\chapter{جریان کاری کامل}
\label{chap:pipeline}

\section{جایگاه در نقشه}

این فصل همان نمودار فصل~\ref{chap:map} را با یک مثال پیمایش می‌کند تا مسیر از فیزیک تا برآورد منابع در ذهن یک خط شود.

\section{ایدهٔ شهودی}

پس از داشتن $H$ دو پرسش جدا می‌شوند: چه چیزی تحول می‌کند (مسئلهٔ دینامیکی: $U(t)=e^{-iHt}$) و چه چیزی خواسته می‌شود (مسئلهٔ ایستا / طیفی). هر پرسش یک یا چند رهیافت دارد؛ هر رهیافت ابزار می‌گیرد؛ ابزار به مدار یا زمان‌بندی سخت‌افزار تبدیل می‌شود؛ خوانش یک مشاهده‌پذیر می‌دهد؛ راستی‌آزمایی کلاسیک می‌گوید عدد باورپذیر است؛ برآورد منابع می‌گوید این مسیر شدنی است یا نه. شکل~\ref{fig:pipeline} همین زنجیره را نشان می‌دهد.

\begin{figure}[htbp]
	\centering
	\begin{latin}
		\setstretch{1}
		\resizebox{\ifdim\width>\textwidth\textwidth\else\width\fi}{!}{%
			% سبک مشترک نمودارهای سند پایه‌ها (فقط داخل محیط latin فراخوانی شود)
\tikzset{
	qafbox/.style={
		draw=black!55,
		rounded corners=2pt,
		align=center,
		inner sep=3pt,
		thick,
		fill=white,
		font=\footnotesize\linespread{1}\selectfont,
		minimum height=0.95cm
	},
	qafwide/.style={
		qafbox,
		text width=2.55cm
	},
	qafnarr/.style={
		qafbox,
		text width=2.05cm
	},
	qaftitle/.style={
		font=\small\bfseries\linespread{1}\selectfont,
		align=center
	},
	qafarr/.style={
		-{Latex[length=2.2mm]},
		thick,
		draw=black!65
	},
	qafpanel/.style={
		draw,
		thick,
		rounded corners=4pt,
		inner xsep=8pt,
		inner ysep=8pt
	}
}

\begin{tikzpicture}[
	font=\scriptsize\linespread{1}\selectfont,
	node distance=0.38cm and 0.38cm
	]
	\node[qaftitle] (t0) {From physical problem to a checked observable};
	\node[qafwide, fill=blue!10, below=0.45cm of t0] (phys) {Physical problem};
	\node[qafwide, fill=blue!10, below=of phys] (model) {Mathematical model};
	\node[qafwide, fill=blue!12, below=of model] (H) {Hamiltonian $H$};

	\node[qafnarr, fill=green!12, below left=0.55cm and 0.2cm of H] (q1) {What evolves?};
	\node[qafnarr, fill=orange!12, below right=0.55cm and 0.2cm of H] (q2) {What is wanted?};
	\node[qafnarr, fill=green!8, below=of q1] (U) {$U(t)=e^{-iHt}$};
	\node[qafnarr, fill=orange!8, below=of q2] (gs) {Ground state /\\spectrum};

	\node[qafbox, text width=1.45cm, fill=green!6, below left=0.45cm and -0.2cm of U] (tr) {Trotter};
	\node[qafbox, text width=1.45cm, fill=green!6, below=0.45cm of U] (lc) {LCU};
	\node[qafbox, text width=1.45cm, fill=green!6, below right=0.45cm and -0.2cm of U] (qs) {QSP};
	\node[qafbox, text width=1.55cm, fill=yellow!12, below left=0.45cm and -0.05cm of gs] (vq) {VQE};
	\node[qafbox, text width=1.55cm, fill=orange!8, below right=0.45cm and -0.05cm of gs] (qp) {QPE};

	\node[qafwide, fill=gray!12, below=1.35cm of lc] (alg) {Quantum algorithm};
	\node[qafwide, fill=gray!10, below=of alg] (hw) {Circuit / adiabatic path};
	\node[qafnarr, fill=cyan!10, below left=0.45cm and 0.05cm of hw] (dig) {Digital};
	\node[qafnarr, fill=cyan!10, below right=0.45cm and 0.05cm of hw] (an) {Analog / hybrid};
	\node[qafwide, fill=gray!8, below=1.25cm of hw] (me) {Measurement};
	\node[qafwide, fill=gray!8, below=of me] (ob) {Observable};
	\node[qafwide, fill=gray!8, below=of ob] (va) {Classical validation};
	\node[qafwide, fill=gray!8, below=of va] (re) {Resource estimation};

	\draw[qafarr] (phys) -- (model);
	\draw[qafarr] (model) -- (H);
	\draw[qafarr] (H.south) -- ++(0,-0.12) -| (q1.north);
	\draw[qafarr] (H.south) -- ++(0,-0.12) -| (q2.north);
	\draw[qafarr] (q1) -- (U);
	\draw[qafarr] (q2) -- (gs);
	\draw[qafarr] (U) -- (tr);
	\draw[qafarr] (U) -- (lc);
	\draw[qafarr] (U) -- (qs);
	\draw[qafarr] (gs) -- (vq);
	\draw[qafarr] (gs) -- (qp);
	\draw[qafarr] (tr.south) |- (alg.west);
	\draw[qafarr] (qs.south) |- (alg.east);
	\draw[qafarr] (lc) -- (alg);
	\draw[qafarr] (vq.south) |- (alg.east);
	\draw[qafarr] (qp.south) |- (alg.east);
	\draw[qafarr] (alg) -- (hw);
	\draw[qafarr] (hw) -- (dig);
	\draw[qafarr] (hw) -- (an);
	\draw[qafarr] (dig.south) |- (me.west);
	\draw[qafarr] (an.south) |- (me.east);
	\draw[qafarr] (me) -- (ob);
	\draw[qafarr] (ob) -- (va);
	\draw[qafarr] (va) -- (re);
\end{tikzpicture}
%
		}
	\end{latin}
	\caption{جریان کاری از مسئلهٔ فیزیکی تا برآورد منابع. انشعاب میانی همان دو پرسش فصل~\ref{chap:dyn-spec} است؛ پایین نمودار انتخاب بستر اجراست نه انتخاب مسئله.}
	\label{fig:pipeline}
\end{figure}

\section{حداقل فرمالیسم}

\paragraph{مثال پیمایش: مدل هابارد.}
\begin{enumerate}
	\item \textbf{مسئلهٔ فیزیکی:} دینامیک الکترون‌های همبسته پس از یک تغییر ناگهانی هامیلتونی، یا انرژی حالت پایه.
	\item \textbf{مدل:} فرمی-هابارد با $t$ و $U$.
	\item \textbf{هامیلتونی:} همان $H$ فصل~\ref{chap:models}.
	\item \textbf{انشعاب سؤال:} اگر \lr{quench} است، رهیافت تحول زمانی $\to$ \lr{Trotter} یا \lr{QSP} $\to$ مدار دیجیتال $\to$ اندازه‌گیری چگالی یا جریان. اگر انرژی پایه است: \lr{VQE} (نزدیک‌مدت)، یا مسیر آدیاباتیک، یا \lr{QPE} روی $U=e^{-iHt}$ (تحمل‌پذیر در برابر خطا).
	\item \textbf{راستی‌آزمایی:} حد $U=0$، شبکهٔ کوچک با قطری‌سازی، تقارن تعداد ذره.
	\item \textbf{برآورد منابع:} تعداد کیوبیت پس از کدگذاری فرمیونی، عمق مدار، گام‌های تروتر یا درجهٔ چندجمله‌ای \lr{QSP}، دقت مشاهده‌پذیر.
\end{enumerate}
همین اسکلت برای \lr{Heisenberg} یا نظریهٔ پیمانه‌ای شبکه‌ای تکرار می‌شود؛ فقط بلوک «مدل» و «کدگذاری» عوض می‌شود.

\section{ارتباط}

نقشهٔ فشرده در فصل~\ref{chap:map}، مدل‌ها در فصل~\ref{chap:models}، دو پرسش در فصل~\ref{chap:dyn-spec}، دو رهیافت در فصل~\ref{chap:evo-adi}، و بستن حلقه در فصل~\ref{chap:closing} آمده است.

\section{چه وقت استفاده می‌شود}

برای دیدن شمای کلی در یک نگاه، و برای توضیح دادن مسیر یک پروژه به خودتان در چند دقیقه.

\begin{remark}[جای توسعه]
نمودار دوم مخصوص سامانهٔ کوانتومی باز (معادلهٔ مادر، مسیرهای کوانتومی)؛ جعبهٔ جدا برای پاسخ خطی؛ نمونهٔ عددی کوچک ($2\times 2$ هابارد).
\end{remark}

% ============================================================
\part{دو پرسش و دو رهیافت}
% ============================================================

\chapter{دو پرسش از یک هامیلتونی: دینامیک و طیف}
\label{chap:dyn-spec}

\section{جایگاه در نقشه}

اولین انشعاب بعد از $H$: دینامیک در برابر ایستا / طیفی. همهٔ ابزارهای بعدی زیر یکی از این دو پرسش می‌نشینند، یا روی پل میانشان.

\section{ایدهٔ شهودی}

هامیلتونی هم مولد حرکت است و هم تعیین‌کنندهٔ انرژی‌های مجاز. از یک $H$ می‌توان دو خانواده سؤال پرسید که الگوریتم‌هایشان فرق دارد. شکل~\ref{fig:two-questions} این انشعاب را نشان می‌دهد.

\begin{figure}[htbp]
	\centering
	\begin{latin}
		\setstretch{1}
		% سبک مشترک نمودارهای سند پایه‌ها (فقط داخل محیط latin فراخوانی شود)
\tikzset{
	qafbox/.style={
		draw=black!55,
		rounded corners=2pt,
		align=center,
		inner sep=3pt,
		thick,
		fill=white,
		font=\footnotesize\linespread{1}\selectfont,
		minimum height=0.95cm
	},
	qafwide/.style={
		qafbox,
		text width=2.55cm
	},
	qafnarr/.style={
		qafbox,
		text width=2.05cm
	},
	qaftitle/.style={
		font=\small\bfseries\linespread{1}\selectfont,
		align=center
	},
	qafarr/.style={
		-{Latex[length=2.2mm]},
		thick,
		draw=black!65
	},
	qafpanel/.style={
		draw,
		thick,
		rounded corners=4pt,
		inner xsep=8pt,
		inner ysep=8pt
	}
}

\begin{tikzpicture}[
	font=\footnotesize\linespread{1}\selectfont,
	node distance=0.5cm and 1.1cm
	]
	\node[qafbox, text width=3.2cm, fill=blue!12] (H) {Hamiltonian $H$};
	\node[qafbox, text width=3.4cm, fill=green!12, below left=of H] (dyn) {%
		\textbf{Dynamic}\\[2pt]
		$U(t)=e^{-iHt}$\\
		$\langle O(t)\rangle$, quench};
	\node[qafbox, text width=3.4cm, fill=orange!12, below right=of H] (stat) {%
		\textbf{Static / spectral}\\[2pt]
		ground state, gap,\\
		eigenvalues};
	\node[qafbox, text width=3.6cm, fill=yellow!14, below=1.35cm of H] (bridge) {%
		Response / Green function:\\
		spectral in appearance,\\
		often needs $f(H)$ or $U(t)$};
	\draw[qafarr] (H.south) -- ++(0,-0.12) -| (dyn.north);
	\draw[qafarr] (H.south) -- ++(0,-0.12) -| (stat.north);
	\draw[qafarr] (dyn.south) |- (bridge.west);
	\draw[qafarr] (stat.south) |- (bridge.east);
\end{tikzpicture}

	\end{latin}
	\caption{دو پرسش از یک $H$. تابع پاسخ روی مرز است: ظاهراً طیفی است، اما اغلب به $f(H)$ یا تحول زمانی نیاز دارد.}
	\label{fig:two-questions}
\end{figure}

\paragraph{پرسش دینامیکی --- تحول زمانی.}
حالت چطور با زمان عوض می‌شود؟ شکل نوعی: سامانه در $\lvert\psi_0\rangle$ است، $H$ ثابت می‌ماند یا ناگهان عوض می‌شود، و می‌خواهیم $\lvert\psi(t)\rangle$ یا $\langle O(t)\rangle$ را بدانیم. شهود: فیلم حرکت سامانه. تغییر ناگهانی هامیلتونی نمونهٔ استاندارد است.

\paragraph{پرسش ایستا / طیفی.}
چه انرژی‌هایی مجازند؟ حالت پایه چیست؟ گاف چقدر است؟ شهود: طیف خطوط یک اتم، یا کمینهٔ انرژی یک مولکول. اینجا «زمان» هدف نیست؛ ویژه‌مقدارها و ویژه‌بردارها هدف‌اند.

\paragraph{پل: پاسخ خطی.}
تابع گرین یا تابع طیفی ظاهراً طیفی است، اما محاسبهٔ عملی‌اش اغلب به $f(H)$ یا به تحول زمانی نیاز دارد.

\section{حداقل فرمالیسم}

دینامیک بسته (یکانی)، با $\hbar=1$:
\begin{equation}
	\lvert\psi(t)\rangle = U(t)\lvert\psi_0\rangle,\qquad U(t)=e^{-iHt},
\end{equation}
و مشاهده‌پذیر
\begin{equation}
	\langle O(t)\rangle = \langle\psi_0\rvert U^\dagger(t)\,O\,U(t)\lvert\psi_0\rangle.
\end{equation}
ایستا: $H\lvert E_n\rangle = E_n\lvert E_n\rangle$؛ حالت پایه $\lvert E_0\rangle$؛ گاف $\Delta = E_1-E_0$. اگر بتوان $U(t)$ را ساخت، برآورد فاز می‌تواند از ویژه‌فاز $U$ به $E_n$ برسد. پس ابزار دینامیک، ورودی بعضی الگوریتم‌های طیفی است.

\section{ارتباط}

جریان کاری در فصل~\ref{chap:pipeline}، رهیافت‌ها در فصل~\ref{chap:evo-adi}، ساخت $U(t)$ و $f(H)$ در فصل~\ref{chap:digital-tools}، و \lr{QPE} و \lr{VQE} در فصل~\ref{chap:spectral-tools} آمده است.

\section{چه وقت استفاده می‌شود}

اولین سؤال بعد از نوشتن $H$: فیلم می‌خواهم یا طیف؟ تا این را مشخص نکنید، انتخاب بین \lr{Trotter} و \lr{VQE} بی‌معناست.

\begin{remark}[جای توسعه]
دینامیک گرمایی و رسیدن به تعادل گرمایی؛ سامانهٔ کوانتومی باز و معادلهٔ مادر؛ فرمول‌بندی پاسخ خطی با یک مثال.
\end{remark}

\chapter{دو رهیافت: تحول زمانی و محاسبهٔ آدیاباتیک}
\label{chap:evo-adi}

\section{جایگاه در نقشه}

دومین انشعاب بزرگ: \emph{چگونه} از $H$ برای محاسبه استفاده می‌کنیم. ابزارهایی مثل \lr{Trotter} زیرِ رهیافت تحول‌اند؛ مسیر $H(s)$ خودِ رهیافت آدیاباتیک است.

\section{ایدهٔ شهودی}

شکل~\ref{fig:two-paradigms} دو رهیافت و نقطهٔ اتصال‌شان را نشان می‌دهد.

\begin{figure}[htbp]
	\centering
	\begin{latin}
		\setstretch{1}
		\resizebox{\ifdim\width>\textwidth\textwidth\else\width\fi}{!}{%
			% سبک مشترک نمودارهای سند پایه‌ها (فقط داخل محیط latin فراخوانی شود)
\tikzset{
	qafbox/.style={
		draw=black!55,
		rounded corners=2pt,
		align=center,
		inner sep=3pt,
		thick,
		fill=white,
		font=\footnotesize\linespread{1}\selectfont,
		minimum height=0.95cm
	},
	qafwide/.style={
		qafbox,
		text width=2.55cm
	},
	qafnarr/.style={
		qafbox,
		text width=2.05cm
	},
	qaftitle/.style={
		font=\small\bfseries\linespread{1}\selectfont,
		align=center
	},
	qafarr/.style={
		-{Latex[length=2.2mm]},
		thick,
		draw=black!65
	},
	qafpanel/.style={
		draw,
		thick,
		rounded corners=4pt,
		inner xsep=8pt,
		inner ysep=8pt
	}
}

\begin{tikzpicture}[
	font=\footnotesize\linespread{1}\selectfont,
	node distance=0.48cm and 1.2cm
	]
	\node[qafbox, text width=3.6cm, fill=green!12] (evo) {%
		\textbf{Time evolution}\\[2pt]
		implement $U(t)$ or $f(H)$};
	\node[qafbox, text width=3.6cm, fill=orange!12, right=of evo] (adi) {%
		\textbf{Adiabatic computing}\\[2pt]
		the path $H(s)$ \emph{is} the algorithm};
	\node[qafnarr, fill=green!7, below left=0.55cm and 0.05cm of evo] (dig) {Digital:\\product formulas,\\LCU, QSP};
	\node[qafnarr, fill=cyan!10, below right=0.55cm and 0.05cm of evo] (ana) {Analog:\\engineer $H$};
	\node[qafnarr, fill=orange!7, below=0.55cm of adi] (path) {$H(s)=(1-s)H_i+s H_p$};
	\coordinate (mid) at ($(evo)!0.5!(adi)$);
	\node[qafbox, text width=4.4cm, fill=gray!10, below=2.35cm of mid] (join) {%
		Digital adiabatic: Trotterize $H(s_k)$};
	\draw[qafarr] (evo) -- (dig);
	\draw[qafarr] (evo) -- (ana);
	\draw[qafarr] (adi) -- (path);
	\draw[qafarr] (path) -- (join);
	\draw[qafarr] (dig.south) |- (join.west);
\end{tikzpicture}
%
		}
	\end{latin}
	\caption{رهیافت تحول زمانی در برابر محاسبهٔ آدیاباتیک. آنالوگ زیرِ تحول می‌ماند و به گیت تجزیه نمی‌شود؛ مسیر آدیاباتیک را می‌توان با ابزار دیجیتالِ همان ستون چپ گسسته کرد (\lr{digital adiabatic})، پس این دو رقیبِ مطلق نیستند.}
	\label{fig:two-paradigms}
\end{figure}

\paragraph{الف) رهیافت تحول زمانی / شرودینگر.}
هدف: ساختن یا مهندسی $U(t)=e^{-iHt}$ یا به‌طور کلی‌تر یک تابع $f(H)$. دیجیتال: $U(t)$ را به گیت تجزیه می‌کنیم (فصل~\ref{chap:digital-tools}). آنالوگ: خودِ $H$ را روی سخت‌افزار می‌سازیم تا طبیعت انتگرال زمانی را بگیرد (فصل~\ref{chap:hardware}). شهود: سامانه را حرکت می‌دهیم و در زمان $t$ نگاه می‌کنیم.

\paragraph{ب) رهیافت محاسبهٔ آدیاباتیک.}
اینجا محاسبه برابر است با یک مسیر در فضای هامیلتونی‌ها، نه یک $U(t)$ ثابت. معمولاً
\begin{equation}
	H(s) = (1-s)\,H_{\mathrm{init}} + s\,H_{\mathrm{problem}},\qquad s\in[0,1].
\end{equation}
اگر $s(t)$ به‌قدر کافی آرام عوض شود و گاف لحظه‌ای صفر نشود، حالت در پایهٔ لحظه‌ای می‌ماند. شهود: به‌جای شکستن $e^{-iHt}$ به گیت، مسئله را در خودِ هامیلتونی روشن می‌کنید. جزئیات قضیه، پارامترها، و یک مثال کامل \lr{Ising} در بخش‌های~\ref{sec:adiabatic-theorem} و~\ref{sec:adi-ising-example} آمده است.

سه چیز که نباید قاطی شوند:
\begin{itemize}
	\item \textbf{شبیه‌سازی آنالوگ:} $H_{\mathrm{hw}}\approx H_{\mathrm{target}}$؛ لزوماً آدیاباتیک نیست و می‌تواند تحول واقعی زمان را بدهد.
	\item \textbf{محاسبهٔ آدیاباتیک (\lr{AQC}):} مسیر $H(s)$ خودِ الگوریتم است؛ هدف نوعاً حالت پایه است.
	\item \textbf{بازپخت کوانتومی:} خویشاوند نوفه‌دار و بهینه‌سازی‌محور \lr{AQC}؛ با شبیه‌سازی دقیق دینامیک یکی نیست.
\end{itemize}
شکست رهیافت آدیاباتیک همان دینامیک غیرآدیاباتیک است: پرش به ترازهای بالاتر وقتی گاف کوچک یا سرعت زیاد است.

\paragraph{اتصال دو رهیافت.}
مسیر آدیاباتیک را می‌توان با فرمول حاصل‌ضربی دیجیتال کرد (\lr{digital adiabatic}): همان \lr{Trotter}، اما روی $H(s_k)$ در گام‌های گسسته.

\section{قضیهٔ آدیاباتیک: فرمول و پارامترها}
\label{sec:adiabatic-theorem}

فرض کنید خانوادهٔ هامیلتونی‌های
\begin{equation}
	H(s),\qquad s\in[0,1],
\end{equation}
در هر $s$ قطری‌پذیر باشد و ویژه‌مقادیرش را به‌صورت صعودی $E_0(s)\le E_1(s)\le\cdots$ بنویسیم. ویژه‌بردار متناظر با پایه را $\lvert\psi_0(s)\rangle$ می‌نامیم (\emph{ویژه‌حالت لحظه‌ای}\footnote{\lr{instantaneous eigenstate}.}). \emph{گاف لحظه‌ای}\footnote{\lr{instantaneous gap}.} میان پایه و اولین تراز برانگیخته
\begin{equation}
	\label{eq:gap}
	\Delta(s) := E_1(s)-E_0(s)
\end{equation}
است و کمینهٔ آن روی مسیر،
\begin{equation}
	\Delta_{\min} := \min_{s\in[0,1]}\Delta(s),
\end{equation}
پارامتر تعیین‌کنندهٔ سختی مسیر است (شکل~\ref{fig:adiabatic-gap}).

\begin{figure}[htbp]
	\centering
	\begin{latin}
		\setstretch{1}
		\resizebox{0.92\textwidth}{!}{%
			% طیف لحظه‌ای مسیر آدیاباتیک: پایه، گاف، و پرهیز از تقاطع
% سبک مشترک نمودارهای سند پایه‌ها (فقط داخل محیط latin فراخوانی شود)
\tikzset{
	qafbox/.style={
		draw=black!55,
		rounded corners=2pt,
		align=center,
		inner sep=3pt,
		thick,
		fill=white,
		font=\footnotesize\linespread{1}\selectfont,
		minimum height=0.95cm
	},
	qafwide/.style={
		qafbox,
		text width=2.55cm
	},
	qafnarr/.style={
		qafbox,
		text width=2.05cm
	},
	qaftitle/.style={
		font=\small\bfseries\linespread{1}\selectfont,
		align=center
	},
	qafarr/.style={
		-{Latex[length=2.2mm]},
		thick,
		draw=black!65
	},
	qafpanel/.style={
		draw,
		thick,
		rounded corners=4pt,
		inner xsep=8pt,
		inner ysep=8pt
	}
}

\begin{tikzpicture}[
	font=\footnotesize\linespread{1}\selectfont,
	x=1.15cm,
	y=0.85cm
	]
	% محورها
	\draw[->, thick] (0,0) -- (6.4,0) node[right] {$s$};
	\draw[->, thick] (0,-2.2) -- (0,2.6) node[above] {$E(s)$};
	\node[below] at (0,0) {$0$};
	\node[below] at (6,0) {$1$};

	% تراز پایه (پایین) و اولین برانگیخته
	\draw[thick, blue!70!black, smooth]
		plot coordinates {(0,-1.6) (1.5,-1.35) (3,-0.55) (4.5,-1.2) (6,-1.7)};
	\draw[thick, orange!75!black, smooth]
		plot coordinates {(0,1.6) (1.5,1.2) (3,0.35) (4.5,1.15) (6,1.55)};

	% گاف کمینه
	\draw[<->, densely dashed, thick] (3,-0.55) -- (3,0.35)
		node[midway, right, font=\scriptsize, text=black!70] {$\Delta_{\min}$};

	% برچسب‌ها
	\node[blue!70!black, font=\scriptsize, align=left] at (0.85,-2.0) {$E_0(0)$};
	\node[blue!70!black, font=\scriptsize, align=left] at (5.55,-2.05) {$E_0(1)$};
	\node[orange!75!black, font=\scriptsize] at (0.9,2.05) {$E_1(0)$};
	\node[orange!75!black, font=\scriptsize] at (5.5,2.0) {$E_1(1)$};

	% فلش مسیر حالت
	\draw[qafarr, blue!50!black] (0.4,-1.55) -- (5.5,-1.65);
	\node[font=\scriptsize, text=black!55] at (3.2,-1.95) {follow ground state};

	% یادداشت نقاط انتهایی
	\node[font=\scriptsize, align=center, text=black!60] at (0.95, -2.85)
		{$H_B$: easy\\ground state};
	\node[font=\scriptsize, align=center, text=black!60] at (5.2, -2.85)
		{$H_P$: Ising\\ground state};
\end{tikzpicture}
%
		}
	\end{latin}
	\caption{طیف لحظه‌ای در طول مسیر $s$. حالت سامانه باید روی شاخهٔ پایه بماند؛ تنگ‌ترین گاف $\Delta_{\min}$ تعیین می‌کند مسیر چقدر باید آرام باشد.}
	\label{fig:adiabatic-gap}
\end{figure}

\emph{زمان‌بندی}\footnote{\lr{schedule}.} نگاشت زمان فیزیکی به پارامتر بی‌بعد است: $s=s(t)$ با $s(0)=0$ و $s(T)=1$. ساده‌ترین انتخاب خطی است،
\begin{equation}
	s(t)=\frac{t}{T},\qquad t\in[0,T],
\end{equation}
ولی هر $s(t)$ یکنواختِ هموار قابل‌قبول است. سامانه از معادلهٔ شرودینگرِ زمان‌وابسته پیروی می‌کند:
\begin{equation}
	i\,\partial_t\lvert\Psi(t)\rangle = H\bigl(s(t)\bigr)\,\lvert\Psi(t)\rangle.
\end{equation}

\paragraph{بیان کاری قضیه.}
اگر سامانه در $t=0$ در $\lvert\psi_0(0)\rangle$ باشد، گاف در طول مسیر مثبت بماند ($\Delta_{\min}>0$)، و زمان کل $T$ به‌قدر کافی بزرگ باشد، آن‌گاه در $t=T$ حالت به $\lvert\psi_0(1)\rangle$ نزدیک می‌ماند. یک کران کافیِ رایج (نه لزوماً بهینه) این است~\cite{AlbashLidar2018}:
\begin{equation}
	\label{eq:adiabatic-bound}
	T \;\gtrsim\; \frac{\displaystyle\max_{s}\,\bigl\lVert\partial_s H(s)\bigr\rVert}{\Delta_{\min}^{2}}\,\frac{1}{\varepsilon},
\end{equation}
که $\varepsilon$ خطای مجاز در فاصله از پایهٔ نهایی است و $\lVert\cdot\rVert$ هنجار عملگری مناسب. شهود: صورتِ کسر شدت تغییر هامیلتونی را می‌سنجد؛ مخرج می‌گوید گاف کوچک‌تر $\Rightarrow$ مسیر گران‌تر.

\paragraph{پارامترها، یک‌جا.}
\begin{itemize}
	\item $H_{\mathrm{init}}$ یا $H_B$: هامیلتونی آغازین با پایهٔ \emph{آسان} (آماده یا شناخته‌شده).
	\item $H_{\mathrm{problem}}$ یا $H_P$: هامیلتونی مسئله؛ پایهٔ آن همان پاسخی است که می‌خواهیم.
	\item $s$ یا $s(t)$: پارامتر درون‌یابی / زمان‌بندی.
	\item $\Delta(s)$ و $\Delta_{\min}$: گاف لحظه‌ای و کمینه‌اش؛ گلوگاه منابع.
	\item $T$: زمان کل اجرا؛ باید با~\eqref{eq:adiabatic-bound} سازگار باشد.
	\item $\varepsilon$: خطای مجاز نسبت به $\lvert\psi_0(1)\rangle$.
\end{itemize}
برای درون‌یابی خطی $H(s)=(1-s)H_B+s H_P$ داریم $\partial_s H = H_P-H_B$، پس صورتِ~\eqref{eq:adiabatic-bound} مستقل از $s$ و برابر $\lVert H_P-H_B\rVert$ است.

\paragraph{هزینه در رهیافت تحول (برای مقایسه).}
آنجا هزینه به هنجار جمله‌های $H$، زمان فیزیکی $t$، و دقت $\varepsilon$ بستگی دارد؛ شکل دقیق به ابزار فصل~\ref{chap:digital-tools} می‌ماند. اینجا گلوگاه $\Delta_{\min}$ است، نه لزوماً ژرفای مدار گیتی.

\section{ردهٔ مسائل مناسب آدیاباتیک}

محاسبهٔ آدیاباتیک (\lr{AQC}) وقتی طبیعی است که پرسشِ اصلی \emph{حالت پایهٔ یک $H_P$ سخت} باشد، نه دینامیک در زمان ثابت~\cite{Farhi2000,AlbashLidar2018}:
\begin{enumerate}
	\item \textbf{بهینه‌سازی ترکیبی و اسپین کلاسیک.} بسیاری از مسائل \lr{NP}-سخت (برش بیشینه، پوشش رأس، تخصیص، \ldots) به یک هامیلتونی \lr{Ising} کلاسیک نگاشته می‌شوند:
	\begin{equation}
		H_P = \sum_{i} h_i Z_i + \sum_{i<j} J_{ij} Z_i Z_j.
	\end{equation}
	پیکربندی پایه‌ایِ $H_P$ همان حل مسئله است.
	\item \textbf{آماده‌سازی حالت پایهٔ مدل‌های فیزیکی.} وقتی $H_P$ خودِ مدل فیزیکی است (مثلاً \lr{Ising} کوانتومی، یا تقریبی از یک هامیلتونی مولکولی) و می‌خواهید $\lvert E_0\rangle$ را روی سخت‌افزار بنشانید.
	\item \textbf{سخت‌افزاری که مسیر می‌دهد.} بازپخت کوانتومی و بعضی سکوهای آنالوگ ذاتاً یک $H(s)$ اجرا می‌کنند؛ حتی اگر نوفه‌دار باشند، زبان مسئله همان مسیر آدیاباتیک است.
\end{enumerate}
اگر پرسش شما تحول زمانیِ دقیق $e^{-iH t}$ یا یک $f(H)$ عمومی است، رهیافت تحول (با \lr{Trotter}/\lr{LCU}/\lr{QSP}) مناسب‌تر است. اگر گاف مسیر به‌سمت صفر می‌رود (گذار بحرانی، تبهگنی)، آدیاباتیک ساده شکست می‌خورد یا $T$ نجومی می‌شود.

%==========================
%=========================
\section{مثال کاری: از میدان عرضی به \lr{Ising}}
\label{sec:adi-ising-example}

هدف این بخش آن است که ایدهٔ محاسبهٔ آدیاباتیک را در یک مثال کوچک اما کامل دنبال کنیم؛ به‌گونه‌ای که کمیت‌هایی مانند هامیلتونی آغازین، هامیلتونی مسئله، مسیر در فضای هامیلتونی‌ها، گاف طیفی و زمان اجرا صرفاً نام‌گذاری باقی نمانند. برای این منظور، یک سامانهٔ متشکل از دو اسپین-$1/2$، یا معادل آن دو کیوبیت، در نظر می‌گیریم.

ایدهٔ اصلی بسیار ساده است. ابتدا سامانه را در یک حالت پایهٔ شناخته‌شده و به‌آسانی قابل آماده‌سازی قرار می‌دهیم. سپس هامیلتونی را به‌آرامی از یک هامیلتونی ساده به هامیلتونی مسئله تغییر می‌دهیم. اگر تغییر به اندازهٔ کافی آهسته باشد و گاف انرژی در طول مسیر بیش از حد کوچک نشود، سامانه در حالت پایهٔ لحظه‌ای باقی می‌ماند و در انتهای مسیر، حالت پایهٔ هامیلتونی مسئله را به دست می‌آوریم.

در این مثال، هامیلتونی آغازین از یک میدان مغناطیسی عرضی به دست می‌آید و هامیلتونی نهایی یک مدل \lr{Ising} است. سپس به‌صورت صریح نشان می‌دهیم که چرا گاف در میانهٔ مسیر به یک کمینه می‌رسد و چگونه مقدار
$\Delta_{\min}\simeq0.66$
در
$s_*\simeq0.62$
به دست می‌آید.

\paragraph{گام ۱ --- هامیلتونی آغازین: یک میدان عرضی ساده.}

برای یک اسپین-$1/2$ در یک میدان مغناطیسی، انرژی برهم‌کنش میدان و ممان مغناطیسی از نوع انرژی زیمان است:
\begin{equation}
	H_{\mathrm{Z}}=-\boldsymbol{\mu}\cdot\mathbf{B}.
\end{equation}
اگر میدان در راستای محور $x$ قرار داشته باشد، مؤلفهٔ متناظر اسپین در هامیلتونی ظاهر می‌شود. با جذب ثابت‌های فیزیکی مانند اندازهٔ ممان مغناطیسی در پارامتر میدان، می‌توان هامیلتونی یک اسپین را به‌صورت
\begin{equation}
	H_{\mathrm{Z}}=-h_x X
\end{equation}
نوشت که در آن $X$ همان ماتریس پائولی در راستای $x$ و $h_x>0$ شدت مؤثر میدان عرضی است.

برای دو اسپین مستقل، انرژی کل برابر مجموع انرژی هر دو اسپین است. بنابراین
\begin{equation}
	\label{eq:HB}
	H_B=-h_x(X_1+X_2),
	\qquad h_x>0,
\end{equation}
که در آن
\begin{equation}
	X_1=X\otimes I,
	\qquad
	X_2=I\otimes X.
\end{equation}
به بیان فیزیکی، $X_1$ و $X_2$ نشان می‌دهند که میدان عرضی به‌ترتیب روی اسپین اول و دوم عمل می‌کند.

در پایهٔ محاسباتی،
\begin{equation}
	X=
	\begin{pmatrix}
		0&1\\
		1&0
	\end{pmatrix},
\end{equation}
و ویژه‌حالت‌های آن عبارت‌اند از
\begin{equation}
	\lvert+\rangle
	=
	\frac{\lvert0\rangle+\lvert1\rangle}{\sqrt{2}},
	\qquad
	\lvert-\rangle
	=
	\frac{\lvert0\rangle-\lvert1\rangle}{\sqrt{2}},
\end{equation}
به‌طوری‌که
\begin{equation}
	X\lvert+\rangle=+\lvert+\rangle,
	\qquad
	X\lvert-\rangle=-\lvert-\rangle.
\end{equation}

از آنجا که $X_1$ و $X_2$ روی دو کیوبیت متفاوت عمل می‌کنند،
\begin{equation}
	[X_1,X_2]=0,
\end{equation}
و بنابراین می‌توان ویژه‌حالت‌های $H_B$ را به‌صورت حاصل‌ضرب تانسوری ویژه‌حالت‌های $X$ ساخت. اگر
$a,b\in\{+,-\}$ باشند، داریم
\begin{equation}
	H_B\bigl(\lvert a\rangle\otimes\lvert b\rangle\bigr)
	=
	-h_x(x_a+x_b)
	\bigl(\lvert a\rangle\otimes\lvert b\rangle\bigr),
\end{equation}
که در آن $x_+=+1$ و $x_-=-1$ هستند.

در نتیجه، چهار حالت ویژه و انرژی‌های متناظر عبارت‌اند از
\begin{equation}
	\begin{array}{c|c}
		\text{حالت} & \text{انرژی}\\
		\hline
		\lvert++\rangle & -2h_x\\
		\lvert+-\rangle & 0\\
		\lvert-+\rangle & 0\\
		\lvert--\rangle & +2h_x
	\end{array}
\end{equation}
و بنابراین، برای $h_x>0$ حالت پایه یکتا است:
\begin{equation}
	\label{eq:initial-ground-state}
	\lvert\psi_0(0)\rangle
	=
	\lvert+\rangle\otimes\lvert+\rangle,
	\qquad
	E_0(0)=-2h_x.
\end{equation}

اگر آن را در پایهٔ محاسباتی بنویسیم،
\begin{equation}
	\label{eq:plus-plus}
	\lvert++\rangle
	=
	\frac{1}{2}
	\left(
	\lvert00\rangle+
	\lvert01\rangle+
	\lvert10\rangle+
	\lvert11\rangle
	\right).
\end{equation}
این حالت به‌سادگی با اعمال دو گیت هادامارد روی $\lvert00\rangle$ آماده می‌شود:
\begin{equation}
	\lvert++\rangle
	=
	(H\otimes H)\lvert00\rangle.
\end{equation}

بنابراین هامیلتونی $H_B$ یک ویژگی بسیار مهم دارد: هم تعریف فیزیکی آن ساده است و هم حالت پایهٔ آن را بدون حل یک مسئلهٔ پیچیده می‌دانیم. این همان نقش «نقطهٔ شروع آسان» در محاسبهٔ آدیاباتیک است.

\paragraph{گام ۲ --- هامیلتونی مسئله: مدل \lr{Ising}.}

اکنون هامیلتونی‌ای را تعریف می‌کنیم که حالت پایهٔ آن پاسخ مسئلهٔ مورد نظر را در خود رمزگذاری کند:
\begin{equation}
	\label{eq:HP}
	H_P
	=
	-JZ_1Z_2-h_z(Z_1+Z_2),
	\qquad
	J>0,\quad h_z>0.
\end{equation}

در اینجا دو نوع برهم‌کنش داریم. جملهٔ
\begin{equation}
	-JZ_1Z_2
\end{equation}
یک جفت‌شدگی فرومغناطیسی است؛ بنابراین پیکربندی‌هایی که دو اسپین در یک جهت قرار دارند، یعنی $\lvert00\rangle$ و $\lvert11\rangle$، انرژی کمتری دارند. جملهٔ
\begin{equation}
	-h_z(Z_1+Z_2)
\end{equation}
یک میدان طولی در راستای $z$ است که بین $\lvert00\rangle$ و $\lvert11\rangle$ تمایز ایجاد می‌کند.

در پایهٔ محاسباتی، چون
\begin{equation}
	Z\lvert0\rangle=+\lvert0\rangle,
	\qquad
	Z\lvert1\rangle=-\lvert1\rangle,
\end{equation}
می‌توان انرژی هر چهار پیکربندی را مستقیماً محاسبه کرد:
\begin{equation}
	\begin{aligned}
		E(\lvert00\rangle)
		&=-J-2h_z,\\
		E(\lvert11\rangle)
		&=-J+2h_z,\\
		E(\lvert01\rangle)
		&=+J,\\
		E(\lvert10\rangle)
		&=+J.
	\end{aligned}
\end{equation}

از آنجا که $J>0$ و $h_z>0$،
\begin{equation}
	E(\lvert00\rangle)
	<
	E(\lvert11\rangle)
	<
	E(\lvert01\rangle)=E(\lvert10\rangle)
\end{equation}
و بنابراین حالت پایهٔ هامیلتونی مسئله یکتا است:
\begin{equation}
	\label{eq:problem-ground-state}
	\lvert\psi_0(1)\rangle=\lvert00\rangle,
	\qquad
	E_0(1)=-J-2h_z.
\end{equation}

این مثال کوچک، ایدهٔ نگاشت مسئلهٔ بهینه‌سازی به یک هامیلتونی را نیز نشان می‌دهد. در یک مسئلهٔ بزرگ‌تر، ضرایب $J_{ij}$ و $h_i$ می‌توانند از تابع هزینهٔ یک مسئلهٔ بهینه‌سازی به دست آیند. در این صورت، هر رشتهٔ بیتی
\begin{equation}
	z=(z_1,\ldots,z_n),
	\qquad z_i\in\{+1,-1\},
\end{equation}
یک پیکربندی ممکن است و انرژی متناظر آن نقش تابع هزینه را بازی می‌کند. یافتن حالت پایهٔ $H_P$ در واقع به معنای یافتن پیکربندی با کمترین هزینه است.

برای دو کیوبیت، این مسئله با بررسی چهار حالت به‌سادگی حل می‌شود؛ اما برای $n$ کیوبیت، فضای پیکربندی دارای
\begin{equation}
	2^n
\end{equation}
حالت است. هدف روش آدیاباتیک آن نیست که این $2^n$ حالت را یکی‌یکی جست‌وجو کند، بلکه تلاش می‌کند از تکامل کوانتومی سامانه برای رسیدن مستقیم به حالت پایه استفاده کند.

\paragraph{گام ۳ --- ساخت مسیر آدیاباتیک.}

اکنون دو انتهای مسئله را داریم:
\begin{equation}
	H_B
	\longrightarrow
	H_P.
\end{equation}
برای اتصال این دو، یک درون‌یابی خطی استاندارد در نظر می‌گیریم~\cite{Farhi2000}:
\begin{equation}
	\label{eq:adiabatic-interpolation-example}
	H(s)
	=
	(1-s)H_B+sH_P,
	\qquad
	0\leq s\leq1.
\end{equation}

پارامتر $s$ زمان فیزیکی نیست؛ بلکه یک پارامتر بدون‌بُعد است که مشخص می‌کند در هر لحظه چه سهمی از هامیلتونی آغازین و چه سهمی از هامیلتونی مسئله در سامانه وجود دارد. در یک زمان‌بندی خطی،
\begin{equation}
	s(t)=\frac{t}{T},
	\qquad
	0\leq t\leq T,
\end{equation}
که در آن $T$ کل زمان اجرای فرآیند است.

بنابراین
\begin{equation}
	H(t)
	=
	\left(1-\frac{t}{T}\right)H_B
	+
	\frac{t}{T}H_P.
\end{equation}

برای ادامهٔ محاسبه، پارامترهای عددی زیر را انتخاب می‌کنیم:
\begin{equation}
	\label{eq:ising-example-parameters}
	h_x=1,
	\qquad
	J=1,
	\qquad
	h_z=\frac14.
\end{equation}

در ابتدا،
\begin{equation}
	H(0)=H_B,
\end{equation}
و در انتها،
\begin{equation}
	H(1)=H_P.
\end{equation}
از نتایج دو گام قبل داریم:
\begin{equation}
	\begin{aligned}
		E_0(0)&=-2,
		&
		\lvert\psi_0(0)\rangle&=\lvert++\rangle,
		\\
		E_0(1)&=-1-\frac12=-\frac32,
		&
		\lvert\psi_0(1)\rangle&=\lvert00\rangle.
	\end{aligned}
\end{equation}

نکتهٔ مهم این است که در حالت کلی
\begin{equation}
	[H_B,H_P]\neq0.
\end{equation}
بنابراین حالت پایهٔ $H_B$ و حالت پایهٔ $H_P$ یکسان نیستند و نمی‌توان سامانه را صرفاً با یک تغییر سادهٔ پایه از ابتدا تا انتها توصیف کرد. حالت پایه در طول مسیر نیز تغییر می‌کند.

برای هر مقدار $s$، ویژه‌مقدارهای لحظه‌ای را به‌صورت
\begin{equation}
	H(s)\lvert\psi_k(s)\rangle
	=
	E_k(s)\lvert\psi_k(s)\rangle,
	\qquad
	E_0(s)\leq E_1(s)\leq\cdots
\end{equation}
تعریف می‌کنیم. اگر سامانه از حالت پایه شروع کند و به‌اندازهٔ کافی آهسته تکامل یابد، هدف این است که در طول مسیر در $\lvert\psi_0(s)\rangle$ باقی بماند و در انتها به $\lvert\psi_0(1)\rangle$ برسد.

همچنین
\begin{equation}
	\label{eq:dHds-example}
	\frac{\partial H}{\partial s}
	=
	H_P-H_B.
\end{equation}

در اینجا باید یک تمایز مفهومی مهم را روشن کرد. رابطهٔ
\begin{equation}
	\lvert\psi(t)\rangle
	=
	e^{-iHt}\lvert\psi(0)\rangle
\end{equation}
عملگر تکامل زمانی یک هامیلتونی ثابت را بیان می‌کند؛ خود این عبارت لزوماً «الگوریتم» نیست. در اجرای دیجیتال، الگوریتم‌هایی مانند \lr{Trotter--Suzuki} یا روش‌های مبتنی بر \lr{LCU} و \lr{QSP} برای تقریب یا پیاده‌سازی این تکامل به کار می‌روند. در محاسبهٔ آدیاباتیک نیز ایدهٔ محاسباتی، تغییر آهستهٔ هامیلتونی $H(s)$ و دنبال‌کردن حالت پایهٔ لحظه‌ای است.

\paragraph{گام ۴ --- چرا گاف مهم است؟}

کل استدلال آدیاباتیک حول یک کمیت طیفی مهم می‌چرخد: گاف بین دو تراز انرژی پایین‌تر. آن را به‌صورت
\begin{equation}
	\label{eq:instantaneous-gap-example}
	\Delta(s)
	=
	E_1(s)-E_0(s)
\end{equation}
تعریف می‌کنیم.

در طول مسیر، $\Delta(s)$ ممکن است تغییر کند. کمیتی که برای زمان اجرای آدیاباتیک اهمیت اساسی دارد، کوچک‌ترین مقدار این گاف است:
\begin{equation}
	\label{eq:min-gap-example}
	\Delta_{\min}
	=
	\min_{0\leq s\leq1}\Delta(s).
\end{equation}

علت اهمیت آن را می‌توان به‌صورت شهودی چنین بیان کرد: هرچه دو تراز کم‌انرژی به یکدیگر نزدیک‌تر شوند، جدا نگه‌داشتن دینامیک سامانه در شاخهٔ حالت پایه دشوارتر می‌شود و سامانه به زمان بیشتری برای تکامل نیاز دارد. بنابراین گاف کوچک معمولاً به یک گلوگاه آدیاباتیک تبدیل می‌شود.

برای مثال حاضر، در نقطهٔ شروع داریم
\begin{equation}
	\operatorname{spec}(H_B)=\{-2,0,0,2\},
\end{equation}
پس
\begin{equation}
	\Delta(0)=2.
\end{equation}
در انتهای مسیر نیز
\begin{equation}
	\operatorname{spec}(H_P)
	=
	\left\{
	-\frac32,-\frac12,1,1
	\right\},
\end{equation}
و در نتیجه
\begin{equation}
	\Delta(1)=1.
\end{equation}

اما از این دو مقدار نمی‌توان نتیجه گرفت که کمینهٔ گاف در ابتدا یا انتهای مسیر قرار دارد. ممکن است گاف در میانهٔ مسیر کوچک‌تر شود. در واقع در این مثال دقیقاً همین اتفاق رخ می‌دهد.

\paragraph{گام ۵ --- ماتریس صریح $H(s)$ و منشأ نزدیک‌شدن ترازها.}

برای مشاهدهٔ این موضوع، ابتدا هامیلتونی‌های ابتدا و انتها را در پایهٔ محاسباتی
\begin{equation}
	\{
	\lvert00\rangle,
	\lvert01\rangle,
	\lvert10\rangle,
	\lvert11\rangle
	\}
\end{equation}
می‌نویسیم.

برای $h_x=1$ داریم
\begin{equation}
	H_B=
	\begin{pmatrix}
		0&-1&-1&0\\
		-1&0&0&-1\\
		-1&0&0&-1\\
		0&-1&-1&0
	\end{pmatrix},
\end{equation}
در حالی که
\begin{equation}
	H_P=
	\begin{pmatrix}
		-\frac32&0&0&0\\
		0&1&0&0\\
		0&0&1&0\\
		0&0&0&-\frac12
	\end{pmatrix}.
\end{equation}

بنابراین
\begin{equation}
	\label{eq:Hs-explicit-example}
	H(s)=
	\begin{pmatrix}
		-\frac32s&-(1-s)&-(1-s)&0\\
		-(1-s)&s&0&-(1-s)\\
		-(1-s)&0&s&-(1-s)\\
		0&-(1-s)&-(1-s)&-\frac12s
	\end{pmatrix}.
\end{equation}

این ماتریس نکتهٔ اصلی مثال را بسیار روشن می‌کند. در $s=0$، جمله‌های خارج از قطر ناشی از میدان عرضی $X$ غالب‌اند و حالت پایه تقریباً همان $\lvert++\rangle$ است. با افزایش $s$، جمله‌های قطری ناشی از هامیلتونی \lr{Ising} تقویت می‌شوند و ترجیح سامانه به سمت پیکربندی $\lvert00\rangle$ تغییر می‌کند.

از آنجا که دو هامیلتونی با یکدیگر جابجایی‌پذیر نیستند، ویژه‌حالت‌ها نیز در طول مسیر تغییر می‌کنند. در نتیجه، ویژه‌مقدارها نیز توابع ساده و خطی از $s$ نیستند.

این نکته اهمیت زیادی دارد:
\begin{equation}
	H(s)=(1-s)H_B+sH_P
\end{equation}
در $s$ خطی است، اما به‌طور کلی
\begin{equation}
	E_k(s)\neq(1-s)E_k(0)+sE_k(1).
\end{equation}
ویژه‌مقدارها از قطری‌سازی خود ماتریس $H(s)$ به دست می‌آیند و در نتیجه می‌توانند رفتار غیرخطی داشته باشند. بنابراین محل کمینهٔ گاف نیز لزوماً $s=1/2$ نیست.

\paragraph{گام ۶ --- استفاده از تقارن برای ساده‌تر کردن مسئله.}

برای اینکه منشأ این رفتار را دقیق‌تر ببینیم، می‌توان از تقارن تبادل دو کیوبیت استفاده کرد. دو حالت
\begin{equation}
	\lvert S\rangle
	=
	\frac{\lvert01\rangle+\lvert10\rangle}{\sqrt2},
	\qquad
	\lvert A\rangle
	=
	\frac{\lvert01\rangle-\lvert10\rangle}{\sqrt2}
\end{equation}
را تعریف می‌کنیم که به‌ترتیب حالت متقارن و پادمتقارن هستند.

هامیلتونی $H(s)$ نسبت به تعویض دو کیوبیت متقارن است، بنابراین فضای حالت به زیر‌فضاهای متقارن و پادمتقارن تجزیه می‌شود. به‌طور خاص،
\begin{equation}
	H(s)\lvert A\rangle
	=
	s\lvert A\rangle.
\end{equation}
بنابراین حالت $\lvert A\rangle$ یک ویژه‌حالت با انرژی
\begin{equation}
	E_A(s)=s
\end{equation}
است.

سه حالت
\begin{equation}
	\{
	\lvert00\rangle,
	\lvert S\rangle,
	\lvert11\rangle
	\}
\end{equation}
یک زیر‌فضای متقارن سه‌بعدی می‌سازند. محدودکردن $H(s)$ به این زیر‌فضا، مسئلهٔ $4\times4$ را به ماتریس $3\times3$ زیر کاهش می‌دهد:
\begin{equation}
	\label{eq:Hsym-example}
	H_{\mathrm{sym}}(s)
	=
	\begin{pmatrix}
		-\frac32s&-\sqrt2(1-s)&0\\
		-\sqrt2(1-s)&s&-\sqrt2(1-s)\\
		0&-\sqrt2(1-s)&-\frac12s
	\end{pmatrix}.
\end{equation}

اکنون ویژه‌مقدارهای مربوط به حالت‌های متقارن از حل
\begin{equation}
	\det\!\left[
	H_{\mathrm{sym}}(s)-E I
	\right]=0
\end{equation}
به دست می‌آیند. این معادله یک چندجمله‌ای درجهٔ سوم در $E$ است. برای هدف حاضر، نوشتن جواب بستهٔ آن با رادیکال‌ها کمک چندانی به فهم فیزیک نمی‌کند؛ مهم این است که برای هر $s$، سه ریشهٔ آن را به‌دست آوریم و سپس آنها را با انرژی حالت پادمتقارن $E_A(s)=s$ مقایسه کنیم.

بنابراین روند محاسبهٔ گاف کاملاً مشخص است:
\begin{equation}
	H(s)
	\longrightarrow
	\{E_0(s),E_1(s),E_2(s),E_3(s)\}
	\longrightarrow
	\Delta(s)=E_1(s)-E_0(s).
\end{equation}

\paragraph{گام ۷ --- مشاهدهٔ کمینهٔ گاف.}

اگر $H(s)$ را برای مقادیر مختلف $s$ قطری کنیم و $\Delta(s)$ را محاسبه کنیم، به یک منحنی غیرخطی برای گاف می‌رسیم. برای پارامترهای
\begin{equation}
	h_x=1,\qquad J=1,\qquad h_z=\frac14,
\end{equation}
کمینهٔ گاف در حدود
\begin{equation}
	\label{eq:sstar-example}
	s_*\simeq0.6213
\end{equation}
رخ می‌دهد. در این نقطه،
\begin{equation}
	E_0(s_*)\simeq-1.13550692,
	\qquad
	E_1(s_*)\simeq-0.47664619,
\end{equation}
و بنابراین
\begin{equation}
	\label{eq:gapmin-numerical-example}
	\Delta_{\min}
	=
	E_1(s_*)-E_0(s_*)
	\simeq0.65886072.
\end{equation}
در نتیجه می‌توان با تقریب مناسب نوشت
\begin{equation}
	\boxed{
		s_*\simeq0.62,
		\qquad
		\Delta_{\min}\simeq0.66.
	}
\end{equation}

این نتیجه را می‌توان به‌صورت فیزیکی نیز تفسیر کرد. در ابتدای مسیر، میدان عرضی تعیین‌کنندهٔ ساختار حالت‌های انرژی است. در انتهای مسیر، برهم‌کنش \lr{Ising} و میدان طولی تعیین‌کنندهٔ ساختار انرژی هستند. در ناحیهٔ میانی، این دو اثر با یکدیگر رقابت می‌کنند و دو تراز کم‌انرژی به یکدیگر نزدیک می‌شوند. از آنجا که $h_z\neq0$ است، تبهگنی نهایی میان $\lvert00\rangle$ و $\lvert11\rangle$ شکسته شده و در این مثال با یک تقاطع واقعی در پایین‌ترین ترازها مواجه نیستیم؛ بلکه رفتار به شکل یک \emph{تقاطع اجتنابی} یا \lr{avoided crossing} ظاهر می‌شود.

در نتیجه، مقدار $s_*$ را نمی‌توان صرفاً از روی فرم خطی $H(s)$ حدس زد. باید طیف لحظه‌ای را محاسبه کرد و سپس تابع
\begin{equation}
	\Delta(s)=E_1(s)-E_0(s)
\end{equation}
را روی بازهٔ $0\leq s\leq1$ کمینه کرد.

این مثال کوچک نشان می‌دهد که چرا در یک پروژهٔ آدیاباتیک، صرف نوشتن $H_B$ و $H_P$ کافی نیست. ساختار طیفی مسیر
\begin{equation}
	H_B
	\longrightarrow
	H(s)
	\longrightarrow
	H_P
\end{equation}
خود یک مسئلهٔ محاسباتی مهم است.

\paragraph{گام ۸ --- ارتباط گاف با زمان اجرا.}

اکنون می‌توانیم نقش گاف را در زمان اجرای آدیاباتیک بررسی کنیم. یک بیان ساده‌شده از شرط آدیاباتیک را می‌توان به‌شکل
\begin{equation}
	\label{eq:adiabatic-condition-example}
	T
	\gg
	\frac{
		\displaystyle
		\max_s
		\left|
		\left\langle
		1(s)
		\middle|
		\frac{dH}{ds}
		\middle|
		0(s)
		\right\rangle
		\right|
	}{
		\Delta_{\min}^2
	}
\end{equation}
نوشت، که در آن $\lvert0(s)\rangle$ و $\lvert1(s)\rangle$ به‌ترتیب حالت پایه و اولین حالت برانگیختهٔ لحظه‌ای هستند.

این رابطه تمام جزئیات قضایای دقیق آدیاباتیک را در خود نشان نمی‌دهد و شکل دقیق شرط به زمان‌بندی، مشتقات هامیلتونی و معیار خطا بستگی دارد؛ اما یک پیام بسیار مهم را آشکار می‌کند:
\begin{equation}
	\boxed{
		T\propto\frac{1}{\Delta_{\min}^{\,2}}
	}
\end{equation}


بنابراین گاف کوچک می‌تواند زمان اجرای بسیار بزرگی ایجاد کند.

برای مثال حاضر،
\begin{equation}
	\frac{\partial H}{\partial s}=H_P-H_B
\end{equation}
و هنجار طیفی آن برابر است با
\begin{equation}
	\label{eq:norm-example}
	\left\|H_P-H_B\right\|_2
	\simeq2.369613.
\end{equation}
اگر برای یک برآورد مرتبه‌ای، این مقدار را به‌عنوان مقیاس صورت کسر و
\begin{equation}
	\varepsilon\sim10^{-2}
\end{equation}
را به‌عنوان خطای هدف در نظر بگیریم، خواهیم داشت
\begin{equation}
	T
	\gtrsim
	\frac{2.37}
	{(0.659)^2(10^{-2})}
	\approx
	5.5\times10^2.
\end{equation}
از این رو مقدار
\begin{equation}
	T\sim500
\end{equation}
را می‌توان به‌عنوان یک تخمین مرتبه‌ای برای این مثال کوچک در نظر گرفت، نه به‌عنوان یک کران دقیق و جهان‌شمول.

در اینجا معنای «گاف، گلوگاه است» کاملاً ملموس می‌شود: بیشتر مسیر ممکن است با گاف نسبتاً بزرگ طی شود، اما یک ناحیهٔ باریک در اطراف $s_*\simeq0.62$ وجود دارد که در آن سامانه حساس‌تر است و سرعت تغییر هامیلتونی باید به‌اندازهٔ کافی کم باشد.

\paragraph{گام ۹ --- نقش میدان طولی $h_z$.}

اکنون می‌توانیم اثر انتخاب $h_z=1/4$ را بهتر درک کنیم. اگر
\begin{equation}
	h_z\rightarrow0,
\end{equation}
در انتهای مسیر انرژی دو حالت
$\lvert00\rangle$ و $\lvert11\rangle$
برابر می‌شود:
\begin{equation}
	E(\lvert00\rangle)
	=
	E(\lvert11\rangle)
	=
	-J.
\end{equation}
در نتیجه حالت پایهٔ نهایی دیگر یکتا نیست و تبهگنی در انتهای مسیر ظاهر می‌شود. در چنین شرایطی، گاف مرتبط با مسئله می‌تواند در نزدیکی انتهای مسیر بسیار کوچک شود یا در حد ایده‌آل به صفر برسد.

این مشاهده یک نکتهٔ مهم را نشان می‌دهد: کوچک‌شدن گاف صرفاً یک ویژگی عددی تصادفی نیست، بلکه می‌تواند مستقیماً از ساختار طیفی هامیلتونی مسئله و نحوهٔ شکسته‌شدن یا باقی‌ماندن تبهگنی‌ها ناشی شود.

در مثال حاضر، $h_z>0$ باعث می‌شود
\begin{equation}
	E(\lvert00\rangle)
	<
	E(\lvert11\rangle),
\end{equation}
و در نتیجه حالت پایهٔ نهایی یکتا باقی بماند. با این حال، در میانهٔ مسیر همچنان یک کمینهٔ گاف در
$s\simeq0.62$
وجود دارد.

\paragraph{گام ۱۰ --- اجرا و خوانش.}

اکنون کل فرآیند را می‌توان به‌صورت یک زنجیرهٔ مشخص خلاصه کرد:

\begin{enumerate}
	\item ابتدا دو کیوبیت را آماده می‌کنیم
	\begin{equation}
		\lvert00\rangle
	\end{equation}
	
	
	\item دو گیت هادامارد اعمال می‌کنیم:
	\begin{equation}
		\lvert00\rangle
		\xrightarrow{H\otimes H}
		\lvert++\rangle
		=
		\lvert\psi_0(0)\rangle.
	\end{equation}
	
	\item هامیلتونی را از $H_B$ به $H_P$ تغییر می‌دهیم:
	\begin{equation}
		H(s)
		=
		(1-s)H_B+sH_P,
		\qquad
		s=\frac{t}{T}.
	\end{equation}
	
	\item اگر $T$ به‌اندازهٔ کافی بزرگ باشد، سامانه در تقریب آدیاباتیک در حالت پایهٔ لحظه‌ای باقی می‌ماند:
	\begin{equation}
		\lvert\psi(t)\rangle
		\approx
		e^{i\phi(t)}
		\lvert\psi_0(s(t))\rangle.
	\end{equation}
	عامل $e^{i\phi(t)}$ یک فاز کلی است و در اندازه‌گیری‌های معمولی مشاهده نمی‌شود.
	
	\item در انتهای مسیر،
	\begin{equation}
		s=1,
		\qquad
		H(1)=H_P,
	\end{equation}
	و حالت سامانه باید به حالت پایهٔ $H_P$ نزدیک باشد:
	\begin{equation}
		\lvert\psi(T)\rangle
		\approx
		\lvert00\rangle.
	\end{equation}
	
	\item در پایان هر دو کیوبیت را در پایهٔ $Z$ اندازه‌گیری می‌کنیم. در اجرای ایده‌آل، نتیجهٔ $\lvert00\rangle$ با احتمال نزدیک به یک به دست می‌آید.
\end{enumerate}

بنابراین، مسئلهٔ «یافتن حالت پایهٔ \lr{Ising}» در این روش به یک مسئلهٔ دنبال‌کردن دینامیکی تبدیل می‌شود. ما $H_P$ را مستقیماً برای یافتن حالت پایهٔ آن قطری نمی‌کنیم؛ بلکه سامانه را از یک حالت پایهٔ ساده و شناخته‌شده به‌صورت کنترل‌شده به سمت حالت پایهٔ هامیلتونی مسئله هدایت می‌کنیم.

\paragraph{گام ۱۱ --- تفسیر الگوریتمی و تمایز میان مسئله، هامیلتونی و الگوریتم.}

این مثال برای تفکیک چند مفهوم که گاهی با یکدیگر خلط می‌شوند نیز مفید است. مدل فیزیکی ما یک مدل \lr{Ising} است و هامیلتونی آن $H_P$ است. مسیر آدیاباتیک
\begin{equation}
	H(s)=(1-s)H_B+sH_P
\end{equation}
یک روش محاسباتی برای تبدیل مسئلهٔ یافتن حالت پایه به یک تکامل کوانتومی است. خود
\begin{equation}
	e^{-iHt}
\end{equation}
عملگر تکامل زمانی است، نه نام یک الگوریتم مستقل.

در یک پیاده‌سازی دیجیتال، برای مثال، می‌توان تکامل ناشی از $H(s_k)$ را در نقاط گسستهٔ مسیر تقریب زد و از روش‌هایی مانند \lr{Trotter--Suzuki} استفاده کرد. در مقابل، در یک پیاده‌سازی آنالوگ یا آدیاباتیک، کنترل مستقیم پارامترهای هامیلتونی می‌تواند مسیر $H(s)$ را در سامانهٔ فیزیکی ایجاد کند.

از این رو زنجیرهٔ مفهومی کامل این مثال را می‌توان چنین نوشت:
\begin{equation}
	\boxed{
		\text{مسئلهٔ فیزیکی/بهینه‌سازی}
		\rightarrow
		H_P
		\rightarrow
		H_B
		\rightarrow
		H(s)
		\rightarrow
		\Delta_{\min}
		\rightarrow
		T
		\rightarrow
		\text{تکامل کوانتومی}
		\rightarrow
		\text{اندازه‌گیری}
	}
\end{equation}

این زنجیره یکی از الگوهای اصلی در طراحی یک پروژهٔ شبیه‌سازی یا محاسبهٔ کوانتومی است.

\paragraph{چگونه به سیستم پیچیده‌تر تعمیم می‌یابد.}

برای $n$ اسپین، ساختار کلی مثال بدون تغییر اساسی باقی می‌ماند. یک هامیلتونی آغازین ساده می‌تواند به‌صورت
\begin{equation}
	\label{eq:HB-n}
	H_B=-h_x\sum_{i=1}^{n}X_i
\end{equation}
و هامیلتونی مسئله به‌صورت
\begin{equation}
	\label{eq:HP-n}
	H_P
	=
	\sum_i h_i Z_i
	+
	\sum_{i<j}J_{ij}Z_iZ_j
\end{equation}
تعریف شود. سپس
\begin{equation}
	\label{eq:H-n-interpolation}
	H(s)=(1-s)H_B+sH_P.
\end{equation}

حالت پایهٔ هامیلتونی آغازین همچنان به‌سادگی قابل آماده‌سازی است:
\begin{equation}
	\lvert\psi_0(0)\rangle
	=
	\lvert+\rangle^{\otimes n},
\end{equation}
که با اعمال $n$ گیت هادامارد روی $\lvert0\rangle^{\otimes n}$ ساخته می‌شود.

اما هامیلتونی مسئله در حالت کلی می‌تواند اطلاعات یک مسئلهٔ بهینه‌سازی بسیار دشوار را در خود رمزگذاری کند. در این حالت، تعداد پیکربندی‌های کلاسیکی برابر $2^n$ است و بررسی مستقیم همهٔ آنها با افزایش $n$ به‌سرعت پرهزینه می‌شود.

نکتهٔ مهم این است که افزایش $n$ به‌تنهایی مقدار مشخصی برای گاف تعیین نمی‌کند. گاف تابعی از کل ساختار هامیلتونی و مسیر انتخاب‌شده است. بنابراین باید زنجیرهٔ زیر را در نظر گرفت:
\begin{equation}
	n
	\longrightarrow
	H_n(s)
	\longrightarrow
	\Delta_{\min}(n)
	\longrightarrow
	T(n)
	\longrightarrow
	\text{پیچیدگی محاسباتی}.
\end{equation}

بسته به مسئله، ممکن است گاف به‌صورت چندجمله‌ای کوچک شود،
\begin{equation}
	\Delta_{\min}(n)\sim n^{-\alpha},
\end{equation}
یا در مسائل دشوارتر به‌صورت نمایی کاهش یابد،
\begin{equation}
	\Delta_{\min}(n)\sim e^{-cn}.
\end{equation}
در حالت دوم، با توجه به وابستگی تقریباً
\begin{equation}
	T\propto\frac{1}{\Delta_{\min}^2},
\end{equation}
زمان آدیاباتیک می‌تواند خود به‌صورت نمایی رشد کند. از این رو، صرفاً افزایش تعداد کیوبیت‌ها معیار کافی برای ارزیابی دشواری نیست؛ ساختار طیفی مسیر اهمیت تعیین‌کننده دارد.

در نتیجه، برای یک پروژهٔ واقعی شبیه‌سازی آدیاباتیک، نوشتن $H_P$ تنها نقطهٔ آغاز است. باید مسیر $H(s)$، ساختار ویژه‌مقدارها، کمینهٔ گاف، زمان‌بندی $s(t)$، منابع مورد نیاز برای اجرا و در نهایت احتمال موفقیت در اندازه‌گیری نیز بررسی شوند.

مثال دوکیوبیتی حاضر دقیقاً به این دلیل ارزشمند است که تمام این مراحل را می‌توان به‌صورت کامل و عددی مشاهده کرد:
\begin{equation}
	\boxed{
		\begin{array}{c}
			\text{میدان عرضی ساده}
			\\[2pt]
			\downarrow
			\\[2pt]
			H_B,\quad \lvert++\rangle
			\\[4pt]
			\downarrow
			\\[2pt]
			H(s)=(1-s)H_B+sH_P
			\\[4pt]
			\downarrow
			\\[2pt]
			E_0(s),E_1(s)
			\\[4pt]
			\downarrow
			\\[2pt]
			\Delta(s)=E_1(s)-E_0(s)
			\\[4pt]
			\downarrow
			\\[2pt]
			\Delta_{\min}\simeq0.66
			\quad\text{در}\quad
			s_*\simeq0.62
			\\[4pt]
			\downarrow
			\\[2pt]
			T\sim\mathcal{O}(10^2\!-\!10^3)
			\\[4pt]
			\downarrow
			\\[2pt]
			\lvert00\rangle
			\quad\text{در خوانش نهایی}
		\end{array}
	}
\end{equation}

بنابراین این مثال کوچک، یک نمونهٔ فشرده از منطق کامل محاسبهٔ آدیاباتیک است: یک حالت پایهٔ آسان آماده می‌شود، هامیلتونی در طول یک مسیر کنترل‌شده تغییر می‌کند، گاف لحظه‌ای تعیین می‌کند که کدام ناحیه از مسیر گلوگاه است، و زمان اجرا باید به‌اندازه‌ای بزرگ انتخاب شود که سامانه بتواند این ناحیهٔ حساس را به‌صورت آدیاباتیک پشت سر بگذارد.
%==========================
%==========================

%\section{مثال کاری: از میدان عرضی به \lr{Ising}}
%\label{sec:adi-ising-example}
%
%هدف این بخش یک محاسبهٔ کاملِ کوچک است تا پارامترهای بخش قبل نام‌گذاری صرف نمانند. دو اسپین-$1/2$ (دو کیوبیت) می‌گیریم: آغاز آسان در میدان مغناطیسی عرضی، پایان یک \lr{Ising} که حالت پایه‌اش رمز پاسخ است.
%
%\paragraph{گام ۱ --- هامیلتونی آغازین (آسان).}
%میدان عرضی یکنواخت:
%\begin{equation}
%	\label{eq:HB}
%	H_B = -h_x\,(X_1+X_2),\qquad h_x>0.
%\end{equation}
%در پایهٔ $\lvert\pm\rangle$ از $X$، ویژه‌حالت‌ها ضرب تانسوری‌اند. کمینهٔ انرژی $-2h_x$ است و
%\begin{equation}
%	\lvert\psi_0(0)\rangle = \lvert+\rangle\otimes\lvert+\rangle
%	= \tfrac12\bigl(\lvert00\rangle+\lvert01\rangle+\lvert10\rangle+\lvert11\rangle\bigr).
%\end{equation}
%آماده‌سازی: دو گیت هادامارد روی $\lvert00\rangle$. این همان «میدان مغناطیسی که پایه‌اش را بلدید» است.
%
%\paragraph{گام ۲ --- هامیلتونی مسئله (\lr{Ising}).}
%\begin{equation}
%	\label{eq:HP}
%	H_P = -J\,Z_1 Z_2 - h_z\,(Z_1+Z_2),\qquad J>0,\ h_z>0.
%\end{equation}
%جملهٔ $Z_1 Z_2$ جفت‌شدگی فرومغناطیسی است؛ $h_z$ میدان طولی کوچک برای شکستن تبهگنی $\lvert00\rangle$ و $\lvert11\rangle$. در پایهٔ محاسباتی انرژی‌ها صریح‌اند:
%\begin{equation}
%	\begin{aligned}
%		E(\lvert00\rangle) &= -J-2h_z,\\
%		E(\lvert11\rangle) &= -J+2h_z,\\
%		E(\lvert01\rangle) &= E(\lvert10\rangle) = +J.
%	\end{aligned}
%\end{equation}
%پس برای $J,h_z>0$ حالت پایهٔ یکتا $\lvert00\rangle$ با انرژی $-J-2h_z$ است. شهود نگاشت: اگر $J_{ij}$ و $h_i$ از یک مسئلهٔ بهینه‌سازی بیایند، یافتن این پایه معادل یافتن کمینهٔ کلاسیک است؛ برای $n$ بزرگ، شمارش $2^n$ پیکربندی ناممکن می‌شود. اینجا $n=2$ است تا همه چیز دستی حل شود.
%
%\paragraph{گام ۳ --- مسیر و پارامترهای عددی.}
%درون‌یابی خطی استاندارد محاسبهٔ آدیاباتیک~\cite{Farhi2000}:
%\begin{equation}
%	H(s)=(1-s)\,H_B + s\,H_P.
%\end{equation}
%اعداد نمونه:
%\begin{equation}
%	h_x=1,\qquad J=1,\qquad h_z=\tfrac14.
%\end{equation}
%آن‌گاه
%\begin{equation}
%	\begin{aligned}
%		H(0)&=H_B, & E_0(0)&=-2, & \lvert\psi_0(0)\rangle&=\lvert++\rangle,\\
%		H(1)&=H_P, & E_0(1)&=-1-\tfrac12=-\tfrac32, & \lvert\psi_0(1)\rangle&=\lvert00\rangle.
%	\end{aligned}
%\end{equation}
%و $\partial_s H=H_P-H_B$؛ مقدار عددی هنجارش در گام بعد برای برآورد $T$ می‌آید.
%
%\paragraph{گام ۴ --- گاف و انتخاب $T$.}
%برای $n=2$ ماتریس $4\times4$ی $H(s)$ را می‌توان قطری کرد. با همین پارامترها، کمینهٔ گاف حدود
%\begin{equation}
%	\Delta_{\min}\approx 0.66
%\end{equation}
%نزدیک $s\approx 0.62$ رخ می‌دهد (تقاطع تراز نیست؛ $h_z$ تبهگنی پایانی را برداشته). هنجار $\lVert H_P-H_B\rVert_2\approx 2.37$ است. از~\eqref{eq:adiabatic-bound} با $\varepsilon\sim 10^{-2}$،
%\begin{equation}
%	T \gtrsim \frac{2.37}{(0.66)^{2}\times 10^{-2}} \approx 5\times 10^{2}.
%\end{equation}
%زمان‌بندی خطی $s(t)=t/T$ را می‌گیریم. اگر $h_z\to 0$، در $s=1$ تبهگنی $\lvert00\rangle$/$\lvert11\rangle$ برمی‌گردد و نزدیک پایان $\Delta$ کوچک می‌شود؛ همان‌جا $T$ باید بزرگ شود --- نمونهٔ کوچکِ «گاف، گلوگاه است».
%
%\paragraph{گام ۵ --- اجرا و خوانش.}
%\begin{enumerate}
%	\item حالت اولیه $\lvert++\rangle$ را آماده کنید.
%	\item $H(s(t))$ را از $s=0$ تا $s=1$ در زمان $T$ اعمال کنید (آنالوگ: کنترل میدان‌ها؛ دیجیتال: تروتر روی $H(s_k)$ در گام‌های گسسته).
%	\item در پایان، هر دو کیوبیت را در پایهٔ $Z$ اندازه بگیرید. با احتمال نزدیک به یک، نتیجه $\lvert00\rangle$ است --- همان پایهٔ $H_P$.
%\end{enumerate}
%این یعنی مسئلهٔ «پیدا کردن پایهٔ \lr{Ising}» به یک مسیر کند در فضای هامیلتونی‌ها ترجمه شده است، نه به شمارش پیکربندی‌ها و نه به قطری‌سازی مستقیم $H_P$ برای $n$ بزرگ.
%
%\paragraph{چگونه به سیستم پیچیده تعمیم می‌یابد.}
%برای $n$ اسپین، همان الگو باقی می‌ماند:
%\begin{equation}
%	H_B=-h_x\sum_{i=1}^{n} X_i,\qquad
%	H_P=\sum_i h_i Z_i+\sum_{i<j} J_{ij} Z_i Z_j,
%\end{equation}
%با $H(s)=(1-s)H_B+s H_P$. پایهٔ $H_B$ همیشه $\lvert+\rangle^{\otimes n}$ است و با هادامارد ساخته می‌شود. پایهٔ $H_P$ --- اگر ضرایب از یک مسئلهٔ سخت بیایند --- می‌تواند به‌لحاظ کلاسیک دشوار باشد؛ آدیاباتیک می‌گوید: اگر $\Delta_{\min}$ مسیر بد نباشد و $T$ مطابق~\eqref{eq:adiabatic-bound} انتخاب شود، سامانه خودش را به آن پایه می‌رساند. هزینهٔ واقعی پژوهش، اغلب برآورد یا کنترل همین $\Delta_{\min}$ است، نه نوشتن $H_P$.

\section{ارتباط}

دو پرسش در فصل~\ref{chap:dyn-spec}، ابزار دیجیتال در فصل~\ref{chap:digital-tools}، \lr{QPE} از تحول تغذیه می‌کند و \lr{VQE} رهیافت سومی است (وردشی) در فصل~\ref{chap:spectral-tools}، و آنالوگ در فصل~\ref{chap:hardware} آمده است. مدل‌های اسپینی در فصل~\ref{chap:models} آمده‌اند.

\section{چه وقت استفاده می‌شود}

تحول: دینامیک، پاسخ زمانی، و هر الگوریتمی که $f(H)$ می‌خواهد. آدیاباتیک: آماده‌سازی حالت پایه وقتی گاف مسیر کنترل‌پذیر به نظر می‌رسد، وقتی مسئله به \lr{Ising}/بهینه‌سازی نگاشته می‌شود، یا وقتی سخت‌افزار به‌طور طبیعی یک $H(s)$ می‌دهد. اگر $\Delta_{\min}$ مسیر را نمی‌شناسید، اول همان را بسنجید؛ وگرنه $T$ حدسِ کور است.

\begin{remark}[جای توسعه]
کران‌های دقیق‌تر قضیه (وابستگی به مشتق‌های بالاتر $H(s)$)؛ زمان‌بندی بهینهٔ غیرخطی؛ \lr{QAOA} به‌عنوان گسسته‌سازی الهام‌گرفته از آدیاباتیک؛ عبور از نقطهٔ بحرانی و بستن گاف در حد ترمودینامیکی؛ مثال عددی کامل قطری‌سازی $H(s)$ برای $n=2$ به‌صورت جدول $E_0(s)$، $E_1(s)$.
\end{remark}

% ============================================================
\part{جعبه ابزار}
% ============================================================

\chapter{جعبه ابزار دیجیتال: \lr{Trotter}، \lr{LCU}، \lr{QSP}}
\label{chap:digital-tools}

\section{جایگاه در نقشه}

زیرِ رهیافت تحول زمانی / $f(H)$. این‌ها الگوریتم کامل «حل هابارد» نیستند؛ زیرروال ساختن یک تابع از هامیلتونی‌اند. خروجی‌شان مدار است.

\section{ایدهٔ شهودی}

روی رایانهٔ دیجیتال نمی‌توان $e^{-iHt}$ را به‌صورت یک دکمه داشت. باید آن را از گیت‌های ساده ساخت. سه نسل شهودی: \lr{Trotter} هامیلتونی را تکه‌تکه اعمال می‌کند؛ \lr{LCU} اگر $H$ جمع چند یکانی باشد آن جمع را با کنترل و پس‌گزینش می‌سازد؛ \lr{QSP} وقتی یک بلوکِ کدگذاری‌شده از $H$ دارید تقریباً هر چندجمله‌ای مناسب از آن را با دنبالهٔ دوران می‌سازد. هر سه پیاده‌کنندهٔ $f(H)$ هستند. تحول زمانی $f(x)=e^{-itx}$ فقط یک انتخاب است~\cite{Childs2018,Berry2015,Low2017}.

\paragraph{\lr{Trotter} --- فرمول حاصل‌ضربی.}
اگر $H=\sum_k H_k$ و هر $e^{-iH_k\Delta t}$ آسان باشد،
\begin{equation}
	e^{-iH t}\approx \Bigl(\prod_k e^{-iH_k t/r}\Bigr)^r.
\end{equation}
شهود: در زمان کوتاه، ترتیب اعمال تکه‌ها کمتر مهم است. هرچه $r$ بیشتر، خطای تروتر کمتر و مدار عمیق‌تر. ساده، شفاف، و نقطهٔ شروع تقریباً هر شبیه‌سازی دیجیتال~\cite{Childs2018}.

\paragraph{\lr{LCU} --- ترکیب خطی یکانی‌ها.}
اگر عملگر هدف ترکیب خطی چند یکانی باشد، می‌توان با یک ثبات کمکی، یک تبدیل روی ضرایب، و یک عمل کنترل‌شده آن را بلوک‌کدگذاری کرد. با اندازه‌گیری ثبات کمکی (و گاهی تکرار) بلوک درست برداشته می‌شود. شهود: درِ ورود به الگوریتم‌های جدیدتر شبیه‌سازی؛ زبان مشترک برای نشاندن $H$ داخل یک مدار یکانی بزرگ~\cite{Berry2015}.

\paragraph{\lr{QSP} --- پردازش سیگنال کوانتومی.}
وقتی $H$ (یا بلوکی مرتبط با آن) داخل یک یکانی بزرگ نشسته، می‌توان با دوران‌های تنظیم‌شده تقریباً هر چندجمله‌ای مناسب از آن بلوک را پیاده کرد. \lr{QSVT} تعمیم همین ایده به تبدیل‌های تکین است. شهود: یک چارچوب واحد برای $e^{-iHt}$، فیلتر طیفی، و $H^{-1}$~\cite{Low2017}. کیوبیتیزاسیون همسایهٔ طبیعی \lr{QSP} است: پیاده‌سازی ساخت‌یافته از بلوک‌کدگذاری. جزئیاتش جای توسعه است، نه فصل جدا در نسخهٔ اول.

\section{حداقل فرمالیسم}

\begin{table}[htbp]
	\centering
	\footnotesize
	\setstretch{1.12}
	\setlength{\tabcolsep}{2pt}
	\caption{سه ابزار دیجیتال زیر رهیافت $f(H)$.}
	\label{tab:digital-tools}
	\begin{tabularx}{\textwidth}{@{}
		>{\raggedright\arraybackslash}p{2.6cm}
		>{\raggedright\arraybackslash}p{2.4cm}
		>{\raggedright\arraybackslash}X
		>{\raggedright\arraybackslash}X
		>{\raggedright\arraybackslash}p{3.8cm}
		@{}}
		\toprule
		\rowcolor{gray!14}
		\textbf{ابزار} & \textbf{ورودی شهودی} & \textbf{خروجی} & \textbf{هزینهٔ شهودی} & \textbf{نقطهٔ ضعف} \\
		\midrule
		فرمول حاصل‌ضربی
		& $H=\sum H_k$
		& تقریب $U(t)$
		& رشد با $t$ و تعداد جمله‌ها
		& خطای تروتر؛ عمق زیاد
		\\
		\rowcolor{gray!8}
		\lr{LCU}
		& $H=\sum \alpha_i U_i$
		& بلوک‌کدگذاری
		& وابسته به $\sum\lvert\alpha_i\rvert$
		& کیوبیت کمکی؛ پس‌گزینش
		\\
		\lr{QSP}/\lr{QSVT}
		& بلوک‌کدگذاری $H$
		& چندجمله‌ای $f(H)$
		& نوعاً خطی در درجهٔ $f$
		& طراحی فازها؛ پیش‌نیاز بلوک
		\\
		\bottomrule
	\end{tabularx}
\end{table}

هر سه نخست برای دینامیک و $f(H)$اند. \lr{QPE} از خروجی‌شان به‌عنوان $U$ کنترل‌شده استفاده می‌کند.

\section{ارتباط}

رهیافت تحول در فصل~\ref{chap:evo-adi}، مصرف‌کنندهٔ طیفی در فصل~\ref{chap:spectral-tools}، لایهٔ ابزار در فصل~\ref{chap:three-layers}، و برآورد منابع در فصل~\ref{chap:closing} آمده است.

\section{چه وقت استفاده می‌شود}

شروع و شهود، یا $H$ خیلی محلی: \lr{Trotter}. نیاز به زبان بلوک‌کدگذاری: \lr{LCU}. هدف یک $f(H)$ عمومی با مقیاس بهتر در رژیم تحمل‌پذیر در برابر خطا: \lr{QSP}/\lr{QSVT}.

\begin{remark}[جای توسعه]
مرتبه‌های تروتر و کران خطا؛ \lr{qDRIFT}؛ مدار کیوبیتیزاسیون و انتخاب چندجمله‌ای \lr{QSP}.
\end{remark}

\chapter{ابزارهای طیفی و ترکیبی: \lr{QPE} و \lr{VQE}}
\label{chap:spectral-tools}

\section{جایگاه در نقشه}

زیرِ پرسش ایستا / طیفی. \lr{QPE} روی تقاطع با رهیافت تحول می‌نشیند (به $U$ کنترل‌شده نیاز دارد). \lr{VQE} رهیافت وردشی / ترکیبی است و هم‌خانوادهٔ \lr{Trotter} نیست.

\section{ایدهٔ شهودی}

دو راه خیلی متفاوت برای انرژی یا طیف.

\paragraph{\lr{QPE} --- برآورد فاز.}
اگر یکانی $U$ ویژه‌بردار $\lvert\phi\rangle$ با ویژه‌مقدار $e^{2\pi i\varphi}$ داشته باشد، برآورد فاز رقم‌های $\varphi$ را روی کیوبیت‌های کمکی می‌نویسد. با گرفتن $U=e^{-iHt}$، $\varphi$ به ویژه‌مقدار $H$ ترجمه می‌شود. پیش‌نیاز سخت: توان‌های کنترل‌شدهٔ $U$ --- یعنی همان شبیه‌سازی فصل~\ref{chap:digital-tools}. به همین دلیل \lr{QPE} الگوریتم طیفیِ خالص نیست. رژیم نوعی: تحمل‌پذیر در برابر خطا، دقت قابل تنظیم، هزینهٔ زیاد مدار~\cite{NielsenChuang2010,McArdle2020}.

\paragraph{\lr{VQE} --- الگوریتم وردشی.}
یک مدار پارامتری $\lvert\psi(\theta)\rangle$ آماده می‌شود، $\langle H\rangle$ با اندازه‌گیری برآورد می‌گردد، و یک بهینه‌ساز کلاسیک $\theta$ را عوض می‌کند. حلقه ترکیبی است: کوانتوم برای حالت و مشاهده‌پذیر، کلاسیک برای جستجوی پارامتر~\cite{Peruzzo2014,Tilly2022}. شهود: به‌جای خواندن طیف با ساعت دقیق، یک آنزات محدود می‌گذارید و انرژی را پایین می‌کشید. مناسب سخت‌افزار نزدیک‌مدت؛ تضمین همگرایی به پایهٔ واقعی ندارد.

\begin{table}[htbp]
	\centering
	\footnotesize
	\setstretch{1.15}
	\caption{دو رژیم برای پرسش طیفی.}
	\label{tab:qpe-vqe}
	\begin{tabularx}{\textwidth}{@{}>{\raggedright\arraybackslash}p{3.3cm} >{\raggedright\arraybackslash}X >{\raggedright\arraybackslash}X@{}}
		\toprule
		\rowcolor{gray!14}
		& \textbf{\lr{QPE} $+$ شبیه‌سازی $U$} & \textbf{\lr{VQE}} \\
		\midrule
		پرسش
		& ویژه‌مقدار / طیف با دقت کنترل‌شده
		& تقریب حالت پایه
		\\
		\rowcolor{gray!8}
		رهیافت
		& تحول کنترل‌شده $+$ برآورد فاز
		& وردشی ترکیبی
		\\
		پیش‌نیاز کوانتومی
		& $U^{2^k}$ کنترل‌شده
		& آنزات کم‌عمق
		\\
		\rowcolor{gray!8}
		رژیم نوعی
		& تحمل‌پذیر در برابر خطا
		& \lr{NISQ}
		\\
		\bottomrule
	\end{tabularx}
\end{table}

سرکوب خطا و تصحیح خطا دو لایهٔ جدا روی این جدول‌اند: اولی غالباً همراه \lr{VQE} و دستگاه نوفه‌دار؛ دومی پیش‌فرض \lr{QPE} مقیاس‌دار~\cite{Preskill2018,Endo2021}.

\section{حداقل فرمالیسم}

\lr{QPE} در یک خط: از $U\lvert\phi\rangle=e^{2\pi i\varphi}\lvert\phi\rangle$ برآورد $\tilde\varphi\approx\varphi$، سپس $E\approx -\varphi/t$ اگر $U=e^{-iHt}$ (با قرارداد فاز مناسب). برای \lr{VQE}
\begin{equation}
	E(\theta)=\langle\psi(\theta)\rvert H\lvert\psi(\theta)\rangle,\qquad
	\theta^\star=\arg\min_\theta E(\theta).
\end{equation}
$H$ به جمله‌های قابل اندازه‌گیری (معمولاً پائولی) شکسته می‌شود و هر کدام جدا خوانش می‌گردد.

\section{ارتباط}

پرسش طیفی در فصل~\ref{chap:dyn-spec}، ساخت $U$ در فصل~\ref{chap:digital-tools}، آدیاباتیک به‌عنوان راه سوم آماده‌سازی پایه در فصل~\ref{chap:evo-adi}، و حلقهٔ ترکیبی در فصل~\ref{chap:hardware} آمده است.

\section{چه وقت استفاده می‌شود}

دقت طیفی در افق \lr{FTQC}: \lr{QPE}. آزمایش روی دستگاه امروز: \lr{VQE}. اگر گاف مسیر معلوم و سخت‌افزار آنالوگ/آدیاباتیک دارید: اول فصل~\ref{chap:evo-adi}.

\begin{remark}[جای توسعه]
برآورد فاز تکیه‌گاه‌شده و کیوبیتیزاسیون$+$\lr{QPE}؛ خانوادهٔ وردشی (\lr{QAOA}، \lr{VQD})؛ نقش تعداد شات در واریانس $\langle H\rangle$.
\end{remark}

% ============================================================
\part{اجرا، راستی‌آزمایی و قطب‌نما}
% ============================================================

\chapter{دیجیتال، آنالوگ، ترکیبی}
\label{chap:hardware}

\section{جایگاه در نقشه}

پایین نمودار، بعد از «الگوریتم»: مدار، مسیر آدیاباتیک، یا زمان‌بندی آنالوگ باید روی یک دستگاه بنشیند. این فصل مدل سخت‌افزاری است، نه فهرست شرکت‌ها.

\section{ایدهٔ شهودی}

\paragraph{دیجیتال --- مدار گیتی.}
الگوریتم به دنبالهٔ گیت از یک مجموعه گیت جهانی ترجمه می‌شود. انعطاف زیاد است؛ هزینه به عمق مدار، همبندی، و خطای گیت حساس است. \lr{Trotter}، \lr{LCU}، \lr{QSP}، \lr{QPE} و \lr{VQE} همگی در نهایت اینجا مدار می‌شوند --- \lr{VQE} معمولاً مدار کم‌عمق‌تر~\cite{Preskill2018}.

\paragraph{آنالوگ --- مهندسی هامیلتونی.}
به‌جای شکستن $U(t)$، $H_{\mathrm{hw}}$ را نزدیک $H_{\mathrm{target}}$ می‌سازید و زمان را واقعی جلو می‌برید. مزیت: عمق گیت معنا ندارد. محدودیت: فقط مدل‌هایی که دستگاه بلداست. یون به‌دام‌افتاده و اتم خنثی، و همچنین آرایه‌های ابررسانا در بعضی رژیم‌ها، می‌توانند چنین کوکی داشته باشند. آنالوگ لزوماً آدیاباتیک نیست؛ می‌تواند همان تحول زمانی مدل هدف باشد~\cite{CiracZoller2012}.

\paragraph{ترکیبی --- حلقه با کلاسیک.}
بخشی از کار روی پردازندهٔ کوانتومی است، بخشی روی کلاسیک. \lr{VQE} نمونهٔ بارز است. سرکوب خطا هم اغلب ترکیبی است. سخت‌افزار بازپیکربندی‌پذیر مرز آنالوگ و دیجیتال را نرم می‌کند: همبندی را می‌توان با مسئله عوض کرد.

\section{حداقل فرمالیسم}

سؤال عملی برای هر پروژه:
\begin{enumerate}
	\item خروجی الگوریتم مدار است، زمان‌بندی $H(t)$ است، یا هر دو؟
	\item دستگاه می‌تواند جمله‌های $H$ را مستقیم بسازد یا فقط گیت جهانی دارد؟
	\item حلقهٔ کلاسیک لازم است (وردشی، کالیبراسیون، سرکوب خطا)؟
\end{enumerate}
اگر پاسخ (۲) مثبت باشد، مسیر آنالوگ را جدی بگیرید. اگر مسئله $f(H)$ عمومی و دقت بالا می‌خواهد، مسیر دیجیتال تحمل‌پذیر در برابر خطا را جدی بگیرید.

\section{ارتباط}

سه مدل محاسباتی در فصل~\ref{chap:qmodel}، آدیاباتیک در برابر تحول در فصل~\ref{chap:evo-adi}، \lr{VQE} در فصل~\ref{chap:spectral-tools}، و اندازه‌گیری و برآورد منابع در فصل~\ref{chap:closing} آمده است.

\section{چه وقت استفاده می‌شود}

وقتی الگوریتم روی کاغذ تمام شده و باید بگویید روی چه موجودی اجرا می‌شود. انتخاب سکو نباید قبل از مشخص شدن پرسش و رهیافت باشد.

\begin{remark}[جای توسعه]
جدول کیفی ابررسانا / یون به‌دام‌افتاده / اتم خنثی؛ محدودیت همبندی و هزینهٔ \lr{SWAP}؛ \lr{analog--digital hybrid}.
\end{remark}

\chapter{بستن حلقه: اندازه‌گیری، راستی‌آزمایی، برآورد منابع}
\label{chap:closing}

\section{جایگاه در نقشه}

انتهای جریان کاری. بدون این سه گام، مدار یا مسیر آدیاباتیک فقط یک آزمایش اجراشده است، نه نتیجهٔ پژوهشی قابل دفاع.

\section{ایدهٔ شهودی}

\paragraph{از مدار تا مشاهده‌پذیر.}
خروجی رایانهٔ کوانتومی تابع موج کامل نیست. خوانش در یک پایه بیت‌های کلاسیک می‌دهد. مشاهده‌پذیر $O$ معمولاً به جمله‌های قابل اندازه‌گیری شکسته می‌شود و میانگین شات‌ها $\langle O\rangle$ را می‌سازد. وفاداری
\ft{fidelity}
 با یک حالت هدف در سامانه‌های بزرگ مستقیماً در دسترس نیست.

\paragraph{راستی‌آزمایی کلاسیک.}
قبل از باور به عدد، آن را جایی که کلاسیک هنوز قوی است چک کنید: شبکه یا مولکول خیلی کوچک (قطری‌سازی یا شبیه‌سازی بردار حالت)؛ حدهای دقیقاً حل‌پذیر ($U=0$ در هابارد، میدان صفر در اسپین)؛ تقارن‌ها (تعداد ذره، اسپین، بار پیمانه‌ای)؛ شبکهٔ تانسوری در رژیم درهم‌تنیدگی کم~\cite{Daley2022,Cirac2021}. اگر الگوریتم در این حدها غلط است، در حد سخت هم به آن اعتماد نکنید.

\paragraph{برآورد منابع.}
سؤال پژوهشی «آیا این مسیر شدنی است؟»: چند کیوبیت پس از کدگذاری؟ عمق مدار یا زمان آنالوگ چقدر است؟ برای \lr{Trotter} چند گام تا خطای تروتر زیر $\varepsilon$؟ برای \lr{QSP} درجهٔ چندجمله‌ای؟ برای \lr{VQE} تعداد پارامتر، شات، و جمله‌های $H$؟ نویز دستگاه در برابر دقت هدف: سرکوب خطا کافی است یا تصحیح خطا لازم است؟

بنچمارک\footnote{\lr{benchmark}. معادل فارسی هنوز در جامعه تثبیت نشده است.} در این سند یعنی مقایسهٔ کنترل‌شده با یک مرجع، نه عدد تبلیغاتی.

\section{حداقل فرمالیسم}

زنجیره:
\begin{equation}
	\text{حالت / مدار}
	\ \longrightarrow\
	\text{خوانش}
	\ \longrightarrow\
	\widehat{\langle O\rangle}
	\ \longrightarrow\
	\text{مرجع کلاسیک}
	\ \longrightarrow\
	\text{هزینه }(n,d,t,\varepsilon).
\end{equation}
برون‌یابی به نویز صفر و لغو احتمالاتی خطا دو نمونه از سرکوب خطا هستند؛ یادگیری نویز پیش‌نیاز بعضی از آن‌هاست~\cite{Endo2021}. هدف نسخهٔ اول فقط این است که این‌ها \emph{بعد از} اندازه‌گیری می‌آیند، نه به‌جای الگوریتم.

\section{ارتباط}

جریان کاری در فصل~\ref{chap:pipeline}، رقبا و حدهای کلاسیک در فصل~\ref{chap:classical}، ابزارها در فصل‌های~\ref{chap:digital-tools} و~\ref{chap:spectral-tools}، و قطب‌نما در فصل~\ref{chap:compass} آمده است.

\section{چه وقت استفاده می‌شود}

پایان هر طرح پروژه، و هر بار که مقاله‌ای فقط مدار را نشان می‌دهد بدون مشاهده‌پذیر، مرجع کلاسیک، یا هزینه.

\begin{remark}[جای توسعه]
جدول شاهدهای راستی‌آزمایی برای هابارد و نظریهٔ پیمانه‌ای؛ فرمول‌های asymptotic هزینه؛ \lr{early-FTQC} میان \lr{NISQ} و تصحیح خطای کامل.
\end{remark}

\chapter{قطب‌نمای پژوهش}
\label{chap:compass}

\section{جایگاه در نقشه}

این فصل برای لحظه‌ای است که مسئلهٔ تازه‌ای روی میز است. اگر فقط دو فصل را باز می‌کنید، فصل~\ref{chap:map} و همین فصل کافی‌اند.

\section{ایدهٔ شهودی}

با چند سؤال پشت‌سرهم جای خود را روی نقشه پیدا کنید. ابزار را آخر انتخاب کنید. شکل~\ref{fig:compass} همان درخت را خلاصه می‌کند.

\begin{figure}[htbp]
	\centering
	\begin{latin}
		\setstretch{1}
		\resizebox{\ifdim\width>\textwidth\textwidth\else\width\fi}{!}{%
			% سبک مشترک نمودارهای سند پایه‌ها (فقط داخل محیط latin فراخوانی شود)
\tikzset{
	qafbox/.style={
		draw=black!55,
		rounded corners=2pt,
		align=center,
		inner sep=3pt,
		thick,
		fill=white,
		font=\footnotesize\linespread{1}\selectfont,
		minimum height=0.95cm
	},
	qafwide/.style={
		qafbox,
		text width=2.55cm
	},
	qafnarr/.style={
		qafbox,
		text width=2.05cm
	},
	qaftitle/.style={
		font=\small\bfseries\linespread{1}\selectfont,
		align=center
	},
	qafarr/.style={
		-{Latex[length=2.2mm]},
		thick,
		draw=black!65
	},
	qafpanel/.style={
		draw,
		thick,
		rounded corners=4pt,
		inner xsep=8pt,
		inner ysep=8pt
	}
}

\begin{forest} qsimtree,
	for tree={
		font=\scriptsize\linespread{1}\selectfont,
		l sep=6mm,
		s sep=1.2mm,
		minimum height=5.2mm,
		inner ysep=1pt,
		align=center
	}
	[{Dynamics or spectrum?}, fill=blue!12
		[{Dynamics / $f(H)$}, fill=green!14
			[{Digital}, fill=green!9
				[{Local $H$: Trotter}, fill=green!5]
				[{General $f(H)$:\\LCU then QSP}, fill=green!5]
			]
			[{Analog:\\engineer $H$}, fill=cyan!10]
		]
		[{Ground state / spectrum}, fill=orange!14
			[{NISQ: VQE}, fill=yellow!16]
			[{Gap along path:\\adiabatic}, fill=orange!8]
			[{High precision:\\simulate $U$, then QPE}, fill=orange!8]
		]
	]
\end{forest}
%
		}
	\end{latin}
	\caption{درخت تصمیم پژوهش. از نوع سؤال شروع کنید؛ بستر و ابزار آخر می‌آیند.}
	\label{fig:compass}
\end{figure}

\section{حداقل فرمالیسم}

چک‌لیست تصمیم؛ هر خط به یک فصل می‌رود.
\begin{enumerate}
	\item \textbf{سؤال من دینامیک است یا طیف؟}
	فیلم $\langle O(t)\rangle$ در برابر انرژی / گاف / ویژه‌بردار.
	$\to$ فصل~\ref{chap:dyn-spec}.
	\item \textbf{مدل کدام است و $H$ چه ساختاری دارد؟}
	اسپین، هابارد، مولکول، نظریهٔ پیمانه‌ای شبکه‌ای؛ محلی بودن؛ فرمیونی یا پائولی.
	$\to$ فصل~\ref{chap:models}.
	\item \textbf{رهیافت: تحول، آدیاباتیک، یا وردشی؟}
	آیا $U(t)$ یا $f(H)$ می‌خواهم، یا مسیر $H(s)$، یا آنزات؟
	$\to$ فصل‌های~\ref{chap:evo-adi} و~\ref{chap:three-layers}.
	\item \textbf{اجرا دیجیتال است، آنالوگ، یا ترکیبی؟}
	مدار گیتی در برابر مهندسی $H$ در برابر حلقه با بهینه‌ساز کلاسیک.
	$\to$ فصل~\ref{chap:hardware}.
	\item \textbf{رژیم دستگاه: \lr{NISQ} یا تحمل‌پذیر در برابر خطا؟}
	\lr{NISQ} غالباً \lr{VQE} و سرکوب خطا؛ \lr{FTQC} غالباً \lr{QSP}/\lr{QPE}.
	$\to$ فصل~\ref{chap:spectral-tools}.
	\item \textbf{اگر دیجیتال و تحول: کدام ابزار؟}
	شروع با \lr{Trotter}؛ $f(H)$ عمومی با \lr{LCU} و \lr{QSP}.
	$\to$ فصل~\ref{chap:digital-tools}.
	\item \textbf{چطور می‌فهمم درست است و آیا منابعش معقول است؟}
	حد کلاسیک، تقارن، برآورد کیوبیت و عمق.
	$\to$ فصل~\ref{chap:closing}.
	\item \textbf{رقیب کلاسیک چیست؟}
	اگر شبکهٔ تانسوری یا \lr{QMC} همان رژیم را خوب می‌پوشاند، مزیت کوانتومی را صریح بنویسید.
	$\to$ فصل~\ref{chap:classical}.
\end{enumerate}

\section{ارتباط}

نقشه در فصل~\ref{chap:map} و واژگان در فصل~\ref{chap:glossary}.

\section{چه وقت استفاده می‌شود}

شروع هر زیرپروژه، خواندن هر مقالهٔ روش‌محور، و هر بار که احساس کردید در جزئیات مدار یا فرمول گم شده‌اید.

\begin{remark}[جای توسعه]
شاخهٔ سامانهٔ کوانتومی باز و معادلهٔ مادر؛ شاخهٔ بهینه‌سازی خالص (\lr{QAOA} / بازپخت) جدا از شبیه‌سازی فیزیکی؛ یک صفحهٔ «پروژهٔ جاری» که پاسخ این سؤال‌ها را برای کار فعلی ثابت نگه دارد.
\end{remark}

\chapter{واژه‌نامهٔ مصوب}
\label{chap:glossary}

منبع واحد اصطلاحات این سند. هنگام افزودن مفهوم تازه، نخست سطر آن را اینجا اضافه کنید، سپس در فصل‌ها به کار ببرید. در متن روان، فارسی مصوب می‌آید و شکل انگلیسی در نخستین کاربرد هر فصل داخل پرانتز است.

\begin{remark}
در مواردی که معادل فارسی هنوز در جامعهٔ علمی داخل تثبیت نشده است (مانند
\lr{benchmark}
یا
\lr{early-FTQC})،
شکل انگلیسی در نخستین کاربرد داخل متن آورده می‌شود تا ابهام ایجاد نشود.
\end{remark}

\begin{longtable}{@{}p{8.8cm}p{5.4cm}@{}}
	\caption{معادل‌های فارسی مصوب این سند برای اصطلاحات پرکاربرد حوزهٔ شبیه‌سازی کوانتومی.}
	\label{tab:glossary}
	\\
	\toprule
	\rowcolor{gray!14}
	\textbf{انگلیسی} & \textbf{فارسی مصوب این سند} \\
	\midrule
	\endfirsthead
	\caption[]{معادل‌های فارسی مصوب این سند (ادامه).}\\
	\toprule
	\rowcolor{gray!14}
	\textbf{انگلیسی} & \textbf{فارسی مصوب این سند} \\
	\midrule
	\endhead
	\bottomrule
	\endfoot
	\lr{gate} & گیت \\
	\rowcolor{gray!8}
	\lr{universal gate set} & مجموعه گیت جهانی \\
	\lr{many-body} & بس‌ذره‌ای \\
	\rowcolor{gray!8}
	\lr{real-time dynamics} & تحول زمانی \\
	\lr{phase estimation} & برآورد فاز \\
	\rowcolor{gray!8}
	\lr{master equation} & معادلهٔ مادر \\
	\lr{transport} & ترابرد \\
	\rowcolor{gray!8}
	\lr{noise} & نویز \\
	\lr{decoherence} & ناهمدوسی \\
	\rowcolor{gray!8}
	\lr{open quantum system} & سامانهٔ کوانتومی باز \\
	\lr{fidelity} & وفاداری \\
	\rowcolor{gray!8}
	\lr{thermalization} & رسیدن به تعادل گرمایی \\
	\lr{variational} & وردشی \\
	\rowcolor{gray!8}
	\lr{fault-tolerant} & تحمل‌پذیر در برابر خطا \\
	\lr{error mitigation} & سرکوب خطا \\
	\rowcolor{gray!8}
	\lr{error correction} & تصحیح خطا \\
	\lr{hybrid} & ترکیبی \\
	\rowcolor{gray!8}
	\lr{Trotter error} & خطای تروتر \\
	\lr{product formula} & فرمول حاصل‌ضربی \\
	\rowcolor{gray!8}
	\lr{tensor network} & شبکهٔ تانسوری \\
	\lr{matrix product state} & حالت حاصل‌ضرب ماتریسی \\
	\rowcolor{gray!8}
	\lr{state vector simulation} & شبیه‌سازی بردار حالت \\
	\lr{connectivity} & همبندی \\
	\rowcolor{gray!8}
	\lr{readout} & خوانش \\
	\lr{vibronic} & ارتعاشی--الکترونی \\
	\rowcolor{gray!8}
	\lr{non-adiabatic} & غیرآدیاباتیک \\
	\lr{singlet fission} & شکافت سینگلت \\
	\rowcolor{gray!8}
	\lr{bosonic mode} & مد بوزونی \\
	\lr{active space} & فضای فعال \\
	\rowcolor{gray!8}
	\lr{trapped ion} & یون به‌دام‌افتاده \\
	\lr{neutral atom} & اتم خنثی \\
	\rowcolor{gray!8}
	\lr{reconfigurable} & بازپیکربندی‌پذیر \\
	\lr{scaling law} & قانون مقیاس‌بندی \\
	\lr{workflow} & جریان کاری \\
	\rowcolor{gray!8}
	\lr{benchmark} & بنچمارک \\
	\lr{random circuit} & مدار تصادفی \\
	\rowcolor{gray!8}
	\lr{dual-unitary} & دوگان‌واحد \\
	\lr{Floquet circuit} & مدار فلوکه \\
	\rowcolor{gray!8}
	\lr{measurement-induced transition} & گذار القایی اندازه‌گیری \\
	\lr{scrambling} & پخش اطلاعات \\
	\rowcolor{gray!8}
	\lr{boson sampling} & نمونه‌برداری بوزونی \\
	\lr{lattice gauge theory} & نظریهٔ پیمانه‌ای شبکه‌ای \\
	\rowcolor{gray!8}
	\lr{discrete time crystal} & بلور زمانی گسسته \\
	\lr{verification} & راستی‌آزمایی \\
	\rowcolor{gray!8}
	\lr{zero-noise extrapolation} & برون‌یابی به نویز صفر \\
	\lr{probabilistic error cancellation} & لغو احتمالاتی خطا \\
	\lr{noise learning} & یادگیری نویز \\
	\rowcolor{gray!8}
	\lr{linear combination of unitaries (LCU)} & ترکیب خطی یکانی‌ها \\
	\lr{qubitization} & کیوبیتیزاسیون \\
	\rowcolor{gray!8}
	\lr{quantum signal processing (QSP)} & پردازش سیگنال کوانتومی \\
	\lr{non-Markovian} & غیرمارکوفی \\
	\rowcolor{gray!8}
	\lr{resource estimation} & برآورد منابع \\
	\lr{quench} & تغییر ناگهانی هامیلتونی \\
	\rowcolor{gray!8}
	\lr{quantum trajectories} & مسیرهای کوانتومی \\
	\lr{Stinespring dilation} & تحول اتساع‌یافته \\
	\rowcolor{gray!8}
	\lr{block-encoding} & بلوک‌کدگذاری \\
	\lr{quantum singular value transformation (QSVT)} & تبدیل تکین کوانتومی \\
	\rowcolor{gray!8}
	\lr{QAOA} & \lr{QAOA} \\
	\lr{NISQ} & \lr{NISQ} \\
	\rowcolor{gray!8}
	\lr{early-FTQC} & \lr{early-FTQC} \\
	\lr{quantum annealing} & بازپخت کوانتومی \\
	\rowcolor{gray!8}
	\lr{adiabatic theorem} & قضیهٔ آدیاباتیک \\
	\lr{instantaneous eigenstate} & ویژه‌حالت لحظه‌ای \\
	\rowcolor{gray!8}
	\lr{instantaneous gap} & گاف لحظه‌ای \\
	\lr{schedule} & زمان‌بندی \\
	\rowcolor{gray!8}
	\lr{driver Hamiltonian} & هامیلتونی راننده \\
	\lr{problem Hamiltonian} & هامیلتونی مسئله \\
\end{longtable}
